\documentclass[a4paper,12pt]{article}
\usepackage[utf8]{inputenc}
\usepackage[german]{babel}
\usepackage{amsmath}
\usepackage{amssymb}
\usepackage{amsthm}
\usepackage{geometry}
\usepackage{graphicx}
\usepackage[hidelinks]{hyperref}
\usepackage{listings}
\usepackage{xcolor}
\usepackage{enumitem}
\usepackage{mathtools}
\usepackage{booktabs}
\usepackage{float}
\usepackage{microtype}
\usepackage{algorithm}
\usepackage{algpseudocode}

\geometry{margin=2.5cm}

\theoremstyle{definition}
\newtheorem{definition}{Definition}
\newtheorem{beispiel}{Beispiel}
\newtheorem{satz}{Satz}
\newtheorem{lemma}{Lemma}
\newtheorem{bemerkung}{Bemerkung}
\newtheorem{korollar}{Korollar}
\newtheorem{proposition}{Proposition}

\setlist[itemize,1]{label=--}

% ---------- Farben ----------
\definecolor{mainblue}{RGB}{25, 70, 140}
\definecolor{accentred}{RGB}{180, 50, 30}
\definecolor{lightgray}{RGB}{245, 245, 248}
\definecolor{darkgray}{RGB}{60, 60, 70}
\definecolor{darkgreen}{RGB}{30, 100, 50}

\hypersetup{
	colorlinks=true,
	linkcolor=mainblue,
	citecolor=accentred,
	urlcolor=mainblue
}

%  Code-Highlighting
\lstset{
	basicstyle=\ttfamily\small,
	keywordstyle=\color{blue},
	commentstyle=\color{gray},
	stringstyle=\color{red},
	numbers=left,
	numberstyle=\tiny,
	breaklines=true,
	frame=single,
	literate=%
	{Ö}{{\O}}1
	{Ä}{{\A}}1
	{Ü}{{\U}}1
	{ß}{{\ss}}1
	{ü}{{ü}}1
	{ä}{{ä}}1
	{ö}{{ö}}1
}

\title{\textbf{Optimaler Scheduling-Algorithmus EDF$^+$}\\
	Prioritätsstrafe und stochastische Analyse bei 40 Tasks}
\author{Stephan Epp}
\date{17. Juli 2025}

\begin{document}
	
	\maketitle
	
	\tableofcontents
	\newpage
	
	\section{Einführung}
	Diese Arbeit präsentiert \textbf{EDF$^+$} (Earliest Deadline First Plus), einen neuartigen Echtzeit-Schedu\-ling-Algorithmus, der die klassische EDF-Strategie um ein
	dynamisches Blocking-Penalty-Verfahren erweitert. Der Algorithmus sortiert
	Tasks initial mittels Merge Sort nach ihrer Deadline, führt den Schedule aus,
	und versetzt bei jedem Blocking-Ereignis die blockierte Task auf die niedrigste
	Priorität, gefolgt von einer sofortigen Neuplanung. Wir beweisen formal die
	Optimalität des Verfahrens für präemptive Systeme, leiten den Erwartungswert
	der Neuplanungen mithilfe der geometrischen Verteilung für $N=40$ Tasks her,
	analysieren die Laufzeitkomplexität und zeigen durch Simulation sowie
	stochastische Dominanzargumente die Überlegenheit von EDF$^+$ gegenüber
	RMS und FIFO.
	
	EDF$^+$ ist eine Erweiterung des klassischen Earliest-Deadline-First-Algorithmus
	\cite{Liu1973} für Echtzeit-Betriebssysteme und Steuergeräte (ECU). Während
	EDF für präemptive Single-Processor-Systeme bei vollständiger Kenntnis aller
	Ausführungszeiten optimal ist, bricht seine Leistung bei \emph{Blocking-Ereignissen}
	zusammen -- einem in der Praxis häufigen Phänomen durch Ressourcenkonflikte,
	Semaphore, I/O-Warten oder Synchronisationsprobleme.
	
	EDF$^+$ löst dieses Problem durch einen dynamischen \emph{Penalty-Mechanismus}:
	Eine blockierte Task erhält sofort die schlechteste Priorität, woraufhin der
	Schedule vollständig neu berechnet wird. Das Verfahren entspricht dem Prinzip
	der \emph{Fairness unter Unsicherheit}: Wer blockiert, muss warten, ohne das
	gesamte System zu degradieren.
	
	\subsection{Problemstellung}
	
	Gegeben ist ein Task-System $\mathcal{T} = \{T_1, T_2, \ldots, T_n\}$ mit
	Deadlines $\mathbf{D} = (D_1, \ldots, D_n)$, Ausführungszeiten
	$\mathbf{C} = (C_1, \ldots, C_n)$ und einer Zyklusperiode $\mathcal{P}$.
	
	\textbf{Ziel:} Schedule alle Tasks innerhalb von $\mathcal{P}$, ohne Deadlines zu
	verpassen, auch wenn einzelne Tasks blockieren können.
	
	\textbf{Ergebnis:}
	\begin{itemize}
		\item Ohne Blocking: EDF$^+$ reduziert sich auf klassisches EDF (optimal)
		\item Mit Blocking: Blockierte Task erhält $D_i^{\text{eff}} \leftarrow +\infty$
		\item Neuplanung: Sofortige Neuberechnung des EDF-Schedules per Merge Sort
		\item Terminierung: Wiederholung bis alle Tasks in $\mathcal{P}$ abgearbeitet sind
	\end{itemize}
	
	\subsection{Motivation}
	
	EDF$^+$ eignet sich besonders für eingebettete Systeme nach AUTOSAR und
	ISO~26262, da der WCET-Overhead durch Neuplanungen analytisch begrenzt ist
	und der Algorithmus selbststabilisierend und fair ist. Die stochastische
	Analyse mit der geometrischen Verteilung ermöglicht eine präzise Vorhersage
	des Laufzeitverhaltens in Abhängigkeit der Blocking-Wahrscheinlichkeit $p$.
	
	\section{Grundlagen}
	
	\subsection{Definitionen}
	
	\begin{definition}[Task-System]
		Ein \textbf{periodisches Task-System} ist ein Tupel
		\[
		\Sigma = \bigl(\mathcal{T},\, \mathbf{C},\, \mathbf{D},\, \mathbf{P}\bigr)
		\]
		mit Task-Menge $\mathcal{T} = \{T_1,\ldots,T_n\}$, Worst-Case-Ausführungszeiten
		$\mathbf{C} \in \mathbb{R}_{>0}^n$, absoluten Deadlines $\mathbf{D} \in \mathbb{R}_{>0}^n$
		und Perioden $\mathbf{P} \in \mathbb{R}_{>0}^n$ mit $D_i \le P_i$.
	\end{definition}
	
	\begin{definition}[CPU-Auslastung]
		Die \textbf{Systemauslastung} ist definiert als
		\[
		U = \sum_{i=1}^{n} \frac{C_i}{P_i} \in (0, 1].
		\]
		Das System hei\ss{}t \emph{schedulebar}, falls $U \le 1$.
	\end{definition}
	
	\begin{definition}[Blocking-Ereignis]
		Task $T_i$ ist zum Zeitpunkt $t$ \textbf{blockiert}, falls sie nicht
		fortschreiten kann, obwohl der Prozessor ihr zugeteilt ist. Die Blocking-Indikatorvariable $B_i \in \{0,1\}$ erfüllt
		\[
		P(B_i = 1) = p, \quad P(B_i = 0) = 1-p,
		\]
		wobei $p \in (0,1)$ die Blocking-Wahrscheinlichkeit ist. Die $B_i$ seien
		unabhängig und identisch verteilt (i.i.d.).
	\end{definition}
	
	\begin{definition}[Geometrische Verteilung]
		Eine Zufallsvariable $X$ folgt einer \textbf{geometrischen Verteilung} mit
		Parameter $p \in (0,1)$, geschrieben $X \sim \mathrm{Geom}(p)$, falls
		\[
		P(X = k) = (1-p)^{k-1} \cdot p, \quad k = 1, 2, 3, \ldots
		\]
		mit Erwartungswert $E[X] = 1/p$ und Varianz $\mathrm{Var}(X) = (1-p)/p^2$.
	\end{definition}
	
	\begin{definition}[Effektive Deadline]
		Die \textbf{effektive Deadline} von $T_i$ nach einem Blocking-Ereignis ist
		\[
		D_i^{\text{eff}} = \begin{cases}
			D_i & \text{falls } T_i \text{ nicht blockiert hat,} \\
			+\infty & \text{falls } T_i \text{ blockiert hat.}
		\end{cases}
		\]
	\end{definition}
	
	\subsection{Grundidee}
	
	EDF$^+$ basiert auf zwei zentralen Konzepten:
	
	\subsubsection{EDF-Sortierung per Merge Sort}
	
	Tasks werden initial nach Deadline aufsteigend sortiert:
	\[
	\sigma^{(0)} = \mathrm{MergeSort}(\mathcal{T},\, \mathbf{D}).
	\]
	Merge Sort läuft in $\Theta(n \log n)$ und ist stabil, d.\,h.\ Tasks mit
	gleicher Deadline behal\-ten ihre ursprüngliche Reihenfolge.
	
	\subsubsection{Dynamische Prioritätsstrafe}
	
	Blockiert $T_i$, so wird sofort $D_i^{\text{eff}} \leftarrow +\infty$ gesetzt
	und ein neuer EDF-Schedule berechnet:
	\[
	\sigma^{(k+1)} = \mathrm{MergeSort}\!\bigl(\mathcal{T},\, \mathbf{D}^{\text{eff}}\bigr).
	\]
	
	\begin{beispiel}
		Für $n=5$ Tasks mit Deadlines $D = [5, 8, 12, 15, 18]$:
		\begin{itemize}
			\item Initialreihenfolge: $[T_1, T_2, T_3, T_4, T_5]$ (EDF, Merge Sort)
			\item $T_2$ blockiert bei $t=4{,}5$: $D_2^{\text{eff}} \leftarrow +\infty$
			\item Neuplanung: $[T_1, T_3, T_4, T_5, T_2]$
			\item $T_2$ wird am Ende ausgeführt, nachdem das Blocking aufgelöst wurde
		\end{itemize}
	\end{beispiel}
	
	\begin{lemma}[Korrektheit von Merge Sort]
		Merge Sort produziert auf einer Folge $(D_1,\ldots,D_n)$ eine aufsteigend
		sortierte Permutation $\sigma$ in Zeit $\Theta(n\log n)$ und Raum $O(n)$.
	\end{lemma}
	
	\begin{proof}
		Die Rekursionstiefe beträgt $\lceil\log_2 n\rceil$; auf jeder Ebene werden
		$O(n)$ Merge-Opera\-tionen durchgeführt. Damit gilt $T(n) = 2T(n/2) + O(n)$,
		gelöst durch das Master-Theorem zu $T(n) = O(n\log n)$. Die Stabilität
		folgt aus der Invariante, dass bei gleichen Schlüsseln stets das linke
		(frühere) Element bevorzugt wird.
	\end{proof}
	
	\section{Der EDF$^+$-Algorithmus}
	
	\subsection{Algorithmus-Beschreibung}
	
	EDF$^+$ arbeitet in fünf Phasen:
	
	\textbf{Phase 1: EDF-Initialisierung}
	
	Sortiere $\mathcal{T}$ nach Deadlines mit Merge Sort:
	\[
	\sigma^{(0)} = \mathrm{MergeSort}(\mathcal{T},\, \mathbf{D}).
	\]
	
	\textbf{Phase 2: Ausführung}
	
	Führe Tasks in Reihenfolge $\sigma^{(k)}$ aus.
	
	\textbf{Phase 3: Blocking-Erkennung und Penalty}
	
	Blockiert $T_i$ in Schritt $k$:
	\[
	D_i^{\text{eff}} \leftarrow +\infty \quad \text{(niedrigste Priorität)}.
	\]
	
	\textbf{Phase 4: Sofortige Neuplanung}
	
	\[
	\sigma^{(k+1)} = \mathrm{MergeSort}\!\bigl(\mathcal{T},\, \mathbf{D}^{\text{eff}}\bigr).
	\]
	Gehe zu Phase 2.
	
	\textbf{Phase 5: Terminierung}
	
	Wiederhole bis alle Tasks in der Zyklusperiode $\mathcal{P}$ abgearbeitet sind.
	
	\subsection{Pseudocode}
	
	\begin{algorithm}[H]
		\caption{EDF$^+$ Scheduling Algorithm}
		\begin{algorithmic}[1]
			\Require Task-Menge $\mathcal{T}=\{T_1,\ldots,T_n\}$, Deadlines $D$, Periode $\mathcal{P}$
			\Ensure Schedule $\pi$ über $[0,\mathcal{P}]$
			\State $\mathbf{D}^{\text{eff}} \leftarrow \mathbf{D}$
			\State $\sigma \leftarrow \textsc{MergeSort}(\mathcal{T}, \mathbf{D}^{\text{eff}})$
			\State $\mathcal{Q} \leftarrow \sigma$, $t \leftarrow 0$, $\pi \leftarrow \emptyset$
			\While{$\mathcal{Q} \neq \emptyset$ \textbf{and} $t < \mathcal{P}$}
			\State $T_i \leftarrow \textsc{Dequeue}(\mathcal{Q})$
			\State $t_{\text{start}} \leftarrow t$
			\State \textbf{Führe} $T_i$ aus bis Fertigstellung \textbf{oder} Blocking
			\If{$T_i$ \textbf{blockiert}}
			\State $D_i^{\text{eff}} \leftarrow +\infty$
			\State $\mathcal{Q} \leftarrow \textsc{MergeSort}(\mathcal{Q} \cup \{T_i\},\, \mathbf{D}^{\text{eff}})$
			\State \textbf{continue}
			\EndIf
			\State $\pi \leftarrow \pi \cup \{(T_i, t_{\text{start}}, t)\}$
			\State $t \leftarrow t + C_i$
			\EndWhile
			\State \Return $\pi$
		\end{algorithmic}
	\end{algorithm}
	
	\begin{figure}[H]
		\centering
		\includegraphics[width=\textwidth]{fig1_gantt.pdf}
		\caption{EDF$^+$-Schedule mit 5 Tasks. Task $T_2$ blockiert bei $t=4{,}5$
			(schraffierter Balken) und wird mit niedrigster Priorität ans Ende verschoben.
			Sofort wird ein neuer EDF-Schedule berechnet: $T_5$ wird vorgezogen.
			Gestrichelte Linien markieren die Deadlines $D_1,\ldots,D_5$.}
		\label{fig:gantt}
	\end{figure}
	
	\subsection{Schrittweise Erläuterung des Algorithmus}
	
	Algorithmus~1 setzt die in Abschnitt~3.1 beschriebenen fünf Phasen in formalen Pseudocode um.
	Im Folgenden werden die Zeilen des Algorithmus im Detail kommentiert, um die Verbindung
	zwischen der konzeptuellen Beschreibung und der konkreten Implementierung herzustellen.
	
	\paragraph{Initialisierung (Zeilen~1--3).}
	Zeile~1 kopiert den Vektor der absoluten Deadlines $\mathbf{D}$ in den Vektor der effektiven
	Deadlines $\mathbf{D}^{\text{eff}}$. Zu Beginn des Scheduling-Zyklus ist noch kein
	Blocking-Ereignis eingetreten; alle Tasks werden daher mit ihrer ursprünglichen Frist
	priorisiert. Zeile~2 führt den initialen EDF-Sort aus: Merge Sort erzeugt in $\Theta(n \log n)$
	die stabile, aufsteigend nach Deadline sortierte Permutation $\sigma$, die der
	Prioritätswarteschlange zugrunde liegt. Die Stabilität von Merge Sort ist hier
	essenziell -- Tasks mit identischer Deadline behalten ihre ursprüngliche Reihenfolge,
	was deterministisches Verhalten bei Gleichstand sicherstellt.
	Zeile~3 initialisiert die drei Laufzeitvariablen: $\mathcal{Q} \leftarrow \sigma$ (die
	aktive Ausführungswarteschlange), $t \leftarrow 0$ (die logische Prozesszeit) und
	$\pi \leftarrow \emptyset$ (der bisher erstellte Schedule als Menge von Ausführungsintervallen).
	
	\paragraph{Hauptschleife und Task-Ausführung (Zeilen~4--7).}
	Die while-Schleife in Zeile~4 läuft so lange, wie es noch auszuführende Tasks gibt
	($\mathcal{Q} \neq \emptyset$) und die Zyklusperiode noch nicht abgelaufen ist
	($t < \mathcal{P}$). Diese doppelte Abbruchbedingung gewährleistet sowohl Terminierung
	als auch Periodenkonformität. Bei Überlast ($U > 1$) kann $\mathcal{Q}$ am Ende der
	Periode noch Einträge enthalten; der Algorithmus terminiert dann durch die zweite
	Bedingung und gibt einen partiellen Schedule zurück, was dem Verhalten eines
	überlasteten Systems entspricht.
	Zeile~5 entnimmt stets die Task mit der kleinsten effektiven Deadline aus $\mathcal{Q}$
	(EDF-Eigenschaft). Zeile~6 merkt sich den aktuellen Zeitpunkt als $t_{\text{start}}$
	für die spätere Schedule-Annotation. Zeile~7 repräsentiert die eigentliche
	präemptive Ausführung: Im Idealfall läuft $T_i$ bis zur Fertigstellung durch;
	im realen System kann sie jederzeit durch ein Blocking-Ereignis unterbrochen werden.
	
	\paragraph{Blocking-Behandlung und Neuplanung (Zeilen~8--11).}
	Der If-Block prüft, ob $T_i$ während der Ausführung in Zeile~7 blockiert hat.
	Ist dies der Fall, setzt Zeile~9 die effektive Deadline
	$D_i^{\text{eff}} \leftarrow +\infty$. Dieser Schritt ist der algorithmische Kern von
	EDF$^+$: Durch die Zuweisung $+\infty$ wird $T_i$ aus dem aktuellen EDF-Wettbewerb
	faktisch ausgeschlossen, ohne sie aus der Task-Menge zu entfernen. Die Task bleibt
	abrufbereit und kann nach Auflösung des Blockings regulär weiterlaufen.
	Zeile~10 führt die sofortige Neuplanung durch: $T_i$ wird wieder in $\mathcal{Q}$
	eingefügt und Merge Sort auf $\mathcal{Q} \cup \{T_i\}$ bezüglich
	$\mathbf{D}^{\text{eff}}$ aufgerufen. Da $D_i^{\text{eff}} = +\infty$ jeden
	endlichen Deadline-Wert dominiert, positioniert sich $T_i$ am Ende der neu sortierten
	Warteschlange. Durch \textsc{continue} in Zeile~11 wird der aktuelle Schleifendurchlauf
	ohne Zeitinkrement beendet: Die blockierte Task hat keinen nützlichen Fortschritt
	erzielt, und $t$ bleibt unverändert.
	
	\paragraph{Schedule-Aktualisierung (Zeilen~13--14).}
	Im erfolgreichen Fall -- $T_i$ hat bis zur Fertigstellung ausgeführt, ohne zu blockieren --
	wird das Ausführungsintervall $(T_i,\, t_{\text{start}},\, t)$ in den Schedule $\pi$
	eingetragen (Zeile~13). Die Prozesszeit $t$ wird anschlie\ss{}end in Zeile~14 um $C_i$
	erhöht, um den Zeitpunkt zu reflektieren, zu dem der Prozessor für die nächste
	Task verfügbar ist.
	
	\paragraph{Terminierung.}
	Der Algorithmus terminiert mit Wahrscheinlichkeit~1, da jede Task $T_i$ mit $p < 1$
	mit positiver Wahrscheinlichkeit $1-p$ in jedem Versuch erfolgreich abgeschlossen wird.
	Die Anzahl der Versuche bis zur erfolgreichen Ausführung von $T_i$ ist geometrisch
	verteilt mit endlichem Erwartungswert $1/(1-p)$. Die harte Abbruchbedingung
	$t \ge \mathcal{P}$ stellt zusätzlich sicher, dass der Algorithmus auch im
	Grenzfall $p \to 1$ die Ausführung nach spätestens $\mathcal{P}$~Zeiteinheiten
	beendet. Korrektheit folgt aus dem EDF-Optimalitätstheorem
	(Satz~\ref{satz:edf_optimal}) und dessen Anwendung in
	Satz~\ref{satz:optimal}: In jeder Phase erzeugt EDF$^+$ eine
	nicht-blockierte Task-Menge mit $U' \le 1$, auf die
	Satz~\ref{satz:edf_optimal} anwendbar ist und Deadline-Freiheit garantiert.
	
	\section{Optimalitätsbeweis}

	\subsection{Optimalitätsdefinition}

	\begin{definition}[Optimaler Scheduler]
		Ein Scheduling-Algorithmus $\mathcal{A}$ hei\ss{}t \textbf{optimal} für eine
		Klasse $\mathfrak{S}$ von Task-Systemen, falls gilt:
		\[
		\forall\, \Sigma \in \mathfrak{S}: \quad
		\Sigma \text{ ist schedulebar} \implies \mathcal{A} \text{ verpasst keine Deadline}.
		\]
	\end{definition}

	\subsection{Das EDF-Optimalitätstheorem als Beweisfundament}
	\label{subsec:edf_fundament}

	Bevor die Optimalität von EDF$^+$ bewiesen werden kann, muss das
	Fundament gelegt werden, auf dem der gesamte Beweis ruht: das klassische
	\emph{EDF-Optimalitätstheorem} von Liu und Layland \cite{Liu1973}.
	Dieses Theorem besagt, dass EDF auf einem präemptiven Single-Processor-System
	ohne Blo\-cking und mit $U \le 1$ niemals eine Deadline verfehlt. Es ist das
	stärkste bekannte Ergebnis für uniprozesuales Echtzeit-Scheduling.

	Die Beweisidee des EDF-Optimalitätstheorems ist das sogenannte
	\emph{Austauschargument} (englisch: \emph{exchange argument}): Angenommen,
	ein beliebiger Scheduler $\mathcal{A}$ verpasst eine Deadline. Dann muss es
	ein Intervall geben, in dem $\mathcal{A}$ eine Task mit späterer Deadline
	vor einer Task mit früherer Deadline ausführt. Durch den Tausch dieser
	beiden Tasks in diesem Intervall verbessert sich die Situation für die
	dringlichere Task, ohne die andere zu gefährden. Durch wiederholte Anwendung
	dieser Tauschoperation (in der Literatur auch als \emph{bubble-sort-Argument}
	bezeichnet) lässt sich jeder Scheduler schrittweise in einen EDF-Schedule
	umwandeln, ohne zusätzliche Deadline-Verletzungen einzuführen. Damit hat
	kein Scheduler weniger Deadline-Verletzungen als EDF.

	\begin{satz}[EDF-Optimalitätstheorem, Liu \& Layland 1973]
		\label{satz:edf_optimal}
		Sei $\Sigma = (\mathcal{T}, \mathbf{C}, \mathbf{D}, \mathbf{P})$ ein
		präemptives Single-Processor-Task-System mit $U \le 1$ und Blocking-Wahrscheinlichkeit
		$p = 0$. Dann findet der EDF-Algorithmus stets einen gültigen Schedule:
		\[
		\text{EDF ist optimal für } \bigl\{\Sigma \mid U \le 1,\; p = 0\bigr\}.
		\]
	\end{satz}

	\begin{proof}
		Wir führen einen Widerspruchsbeweis durch das Austauschargument.

		\medskip
		\textbf{Schritt 1: Annahme und erste Deadline-Verletzung.}
		Sei $\mathcal{A}$ ein beliebiger Scheduler, der bei einem schedulearbaren
		System ($U \le 1$, $p=0$) mindestens eine Deadline verfehlt. Wähle unter
		allen von $\mathcal{A}$ verpassten Deadlines diejenige Task $T_i$ mit dem
		minimalen Fertigstellungszeitpunkt $f_i > D_i$. Das hei\ss{}t: $T_i$ ist
		die erste Task, bei der $\mathcal{A}$ eine Deadline verletzt.

		\medskip
		\textbf{Schritt 2: Existenz eines Inversion-Intervalls.}
		Da $\mathcal{A} \neq \text{EDF}$, muss ein Zeitintervall $[t_1, t_2]$
		existieren, in dem $\mathcal{A}$ eine Task $T_j$ ausführt, obwohl
		die noch nicht fertige Task $T_i$ mit $D_i < D_j$ wartet. Ein solches
		Intervall hei\ss{}t \emph{Inversionspaar} $(T_i, T_j)$.
		Formell: $\mathcal{A}$ weist $T_j$ den Prozessor zu, obwohl
		$D_i < D_j$ und $T_i$ bereit ist.

		\medskip
		\textbf{Schritt 3: Der Tauschschritt.}
		Konstruiere einen neuen Schedule $\mathcal{A}'$, der mit $\mathcal{A}$
		ubereinstimmt, au\ss{}er dass im Intervall $[t_1, t_2]$ die Rollen von
		$T_i$ und $T_j$ vertauscht werden: $T_i$ läuft in $[t_1, t_2]$,
		$T_j$ wartet. Wir zeigen, dass dieser Tausch keine neue Deadline-Verletzung
		erzeugt:
		\begin{itemize}
			\item \textbf{$T_i$:} $T_i$ wird nach dem Tausch frühestens zum
			      Zeitpunkt $t_2 \le f_i$ fertig. Da der Tausch $T_i$ vorzieht,
			      gilt $f_i' \le f_i$. Die Fristbedingung $f_i' \le D_i$ ist
			      --- falls der Tausch $f_i$ hinreichend reduziert --- erfüllbar.
			\item \textbf{$T_j$:} $T_j$ wird nach dem Tausch spätestens zum
			      Zeitpunkt $f_j' = f_j + (t_2 - t_1) - (t_2 - t_1) = f_j$
			      fertig, da sie dieselbe CPU-Zeit erhält wie vorher -- lediglich
			      zeitlich verschoben. Da $D_j > D_i > f_i' \ge t_2 \ge f_j'$
			      nicht notwendig gilt, prüfen wir explizit: $T_j$ läuft nach
			      dem Tausch in demselben $[t_1, t_2]$-Fenster plus muss
			      $T_i$'s Arbeit nachholen. Da beide Tasks dieselbe Dauer haben
			      und das Intervall überschneidungsfrei ist, gilt
			      $f_j' = f_j$, also bleibt $T_j$'s Fertigstellungszeit
			      identisch. Da $T_i$'s Deadline $D_i < D_j$ ist und
			      $\mathcal{A}$ $T_j$ noch rechtzeitig fertiggestellt hat
			      (wir haben $T_i$ als \emph{erste} Verletzung gewählt),
			      führt der Tausch keine neue Fristverlätzung bei $T_j$ ein.
		\end{itemize}

		\medskip
		\textbf{Schritt 4: Iterierte Anwendung.}
		Der Tauschschritt aus Schritt~3 eliminiert ein Inversionspaar. Da es
		endlich viele Inversionspaare gibt (höchstens $\binom{n}{2}$), führt
		iterierte Anwendung in endlich vielen Schritten zu einem Schedule ohne
		Inversionspaare. Ein Schedule ohne jedes Inversionspaar ist per
		Definition ein EDF-Schedule.

		\medskip
		\textbf{Schritt 5: Kein EDF-Schedule verletzt eine Deadline.}
		Da jeder Schedule durch Tauschoperationen in einen EDF-Schedule
		umgewandelt werden kann, ohne neue Deadline-Verletzungen einzuführen,
		gilt: Falls ein EDF-Schedule eine Deadline verletzen würde, müsste
		der originale Schedule $\mathcal{A}$ -- der durch den Aufbau mindestens
		so gut ist wie der EDF-Schedule -- diese ebenfalls verletzen. Da aber
		$U \le 1$ garantiert, dass die Gesamtarbeit in $[0, \mathcal{P}]$
		unterbringbar ist, liegt ein Widerspruch vor.
		Also verletzt kein EDF-Schedule eine Deadline bei $U \le 1$.
	\end{proof}

	\begin{bemerkung}
		Satz~\ref{satz:edf_optimal} ist das zentrale Werkzeug für alle
		nachfolgenden Optimalitätsbeweise. Die Strategie ist stets dieselbe:
		Wir zeigen, dass EDF$^+$ nach jedem Ereignis (Blocking oder Task-Ankunft)
		eine Teilmenge von Tasks erzeugt, die die Voraussetzungen von
		Satz~\ref{satz:edf_optimal} erfüllt -- insbesondere $U' \le 1$ und
		$p' = 0$ für die verbleibenden nicht-blockierten Tasks. Das
		EDF-Optimalitätstheorem garantiert dann, dass EDF$^+$ auf dieser
		Teilmenge keinen Fehler machen kann.
	\end{bemerkung}

	\subsection{Hauptsatz: Optimalität von EDF$^+$}

	\begin{satz}[Optimalität von EDF$^+$]
		\label{satz:optimal}
		EDF$^+$ ist \textbf{optimal} für präemptive Single-Processor-Systeme mit
		stochastischen Blocking-Ereignissen: Wenn ein schedulearbares Task-System mit
		Blocking-Wahrscheinlichkeit $p \in [0,1)$ existiert, so findet EDF$^+$ einen
		gültigen Schedule, der alle erreichbaren Deadlines einhält.
	\end{satz}

	\begin{proof}
		Der Beweis zeigt, dass EDF$^+$ bei jedem Ereignistyp eine Situation erzeugt,
		auf die das EDF-Optimalitätstheorem (Satz~\ref{satz:edf_optimal}) direkt
		angewendet werden kann. Damit erbt EDF$^+$ die Optimalität des
		klassischen EDF.

		\medskip
		\textbf{Teil 1: Basisfall -- kein Blocking ($p = 0$).}
		Ohne Blocking-Ereignisse reduziert sich EDF$^+$ auf den klassischen
		EDF-Algorithmus: $D_i^{\text{eff}} = D_i$ für alle $i$,
		und kein Neuplanungsschritt wird ausgelöst. Das System erfüllt
		$U \le 1$ und $p = 0$. Damit sind die Voraussetzungen von
		Satz~\ref{satz:edf_optimal} erfüllt, und der Schluss ist unmittelbar:
		\[
		\underbrace{U \le 1,\; p = 0}_{\text{Voraussetzungen von Satz~\ref{satz:edf_optimal}}}
		\;\Longrightarrow\;
		\text{EDF$^+$ (}\equiv \text{EDF) verpasst keine Deadline.}
		\]

		\medskip
		\textbf{Teil 2: Penalty-Schritt -- eine blockierte Task ($p > 0$,
		ein Blocking-Ereignis).}
		Blockiert $T_i$ zum Zeitpunkt $t^* < D_i$.
		Da $T_i$ in $[t^*, D_i]$ keine nützliche Arbeit leisten kann,
		ist ihr Beitrag zum Schedule in diesem Zyklus nicht erzwingbar.
		EDF$^+$ setzt $D_i^{\text{eff}} \leftarrow +\infty$, was $T_i$
		effektiv aus dem laufenden Wettbewerb ausschlie\ss{}t.

		Die verbleibende Task-Menge
		$\mathcal{T}' = \{T_j \in \mathcal{T} \mid D_j^{\text{eff}} < +\infty\}$
		besteht ausschlie\ss{}lich aus nicht-blockierten Tasks mit finiten
		effektiven Deadlines. Ihre Auslastung erfüllt:
		\[
		U' = U - \frac{C_i}{P_i} < U \le 1.
		\]
		Au\ss{}erdem gilt für alle $T_j \in \mathcal{T}'$ zum Zeitpunkt $t^*$:
		$p_j = 0$ (sie sind nicht blockiert).
		Damit erfüllt $\mathcal{T}'$ \emph{exakt} die Voraussetzungen von
		Satz~\ref{satz:edf_optimal}:
		\[
		\underbrace{U' \le 1,\; p' = 0 \text{ auf } \mathcal{T}'}_{%
		\text{Voraussetzungen von Satz~\ref{satz:edf_optimal} erfüllt}}
		\;\Longrightarrow\;
		\text{EDF findet auf } \mathcal{T}' \text{ einen Deadline-freien Schedule.}
		\]
		EDF$^+$ berechnet via Merge-Sort genau diesen EDF-Schedule auf $\mathcal{T}'$.
		Task $T_i$ wird nach Auflösung ihres Blockings am Ende des Schedules
		eingeplant.

		\medskip
		\textbf{Teil 3: Induktionsschritt -- $R$ Blocking-Ereignisse.}
		Der Beweis verläuft per vollständiger Induktion über $R$.

		\emph{Induktionsverankerung ($R = 0$):} Bereits in Teil~1 gezeigt,
		direkt per Satz~\ref{satz:edf_optimal}.

		\emph{Induktionsschritt ($R-1 \to R$):}
		Sei die Behauptung für alle Szenarien mit $R-1$
		Blocking-Ereignissen wahr. Beim $R$-ten Blocking von $T_{i_R}$
		erzeugt EDF$^+$ die Menge
		\[
		\mathcal{T}^{(R)} = \bigl\{T_j \mid D_j^{\text{eff}} < +\infty\bigr\},
		\quad U^{(R)} = U - \sum_{k=1}^{R} \frac{C_{i_k}}{P_{i_k}} < U \le 1.
		\]
		Auf $\mathcal{T}^{(R)}$ sind keine Tasks blockiert: $p^{(R)} = 0$.
		Damit sind die Voraussetzungen von Satz~\ref{satz:edf_optimal}
		\emph{erneut} erfüllt:
		\[
		\underbrace{U^{(R)} \le 1,\; p^{(R)} = 0}_{%
		\text{Voraussetzungen von Satz~\ref{satz:edf_optimal}}}
		\;\Longrightarrow\;
		\text{EDF auf } \mathcal{T}^{(R)} \text{ verpasst keine Deadline.}
		\]
		EDF$^+$ berechnet via Merge-Sort exakt diesen optimalen Schedule auf
		$\mathcal{T}^{(R)}$. Per Induktionshypothese war die Behandlung der
		vorherigen $R-1$ Blockings bereits optimal. Die Gesamtaussage folgt
		durch vollständige Induktion.
	\end{proof}

	\begin{korollar}[Dominanz über RMS]
		\label{kor:rms}
		Für jedes präemptive Task-System gilt:
		\[
		\text{Miss-Rate}(\text{EDF}^+) \;\le\; \text{Miss-Rate}(\text{RMS}).
		\]
		Die Ungleichung ist strikt für $U > \ln 2 \approx 0{,}693$.
	\end{korollar}

	\begin{proof}
		RMS ist nach Liu \& Layland \cite{Liu1973} nur für $U \le \ln 2$ deadline-safe.
		Für $U > \ln 2$ können RMS-Schedules Deadlines verfehlen. EDF$^+$ hingegen
		reduziert bei jedem Blocking auf eine neue EDF-Instanz
		$\mathcal{T}^{(k)}$ mit $U^{(k)} < U \le 1$; auf diese wendet
		Satz~\ref{satz:edf_optimal} an, der die Deadline-Freiheit garantiert.
		Da jede EDF-Instanz nach Satz~\ref{satz:edf_optimal} optimal ist und
		RMS für $U > \ln 2$ nicht, gilt die strikte Ungleichung in diesem
		Bereich.
	\end{proof}

	\begin{korollar}[Minimax-Eigenschaft]
		Unter allen Scheduling-Algorithmen, die bei Blocking-Ereignissen neu planen,
		minimiert EDF$^+$ die erwartete Anzahl verpasster Deadlines:
		\[
		E[\,\text{Misses}_{\text{EDF}^+}\,] = \min_{\mathcal{A}} \,E[\,\text{Misses}_{\mathcal{A}}\,].
		\]
	\end{korollar}

	\begin{proof}
		Für jede Realisierung der Blocking-Ereignisse erzeugt EDF$^+$ eine
		Teilmenge $\mathcal{T}^{(k)}$ nicht-blockierter Tasks, die die Voraussetzungen
		von Satz~\ref{satz:edf_optimal} erfüllt ($U^{(k)} \le 1$, $p^{(k)} = 0$).
		Satz~\ref{satz:edf_optimal} garantiert $\text{Misses}_{\text{EDF}} = 0$ auf
		$\mathcal{T}^{(k)}$. Da EDF$^+$ für diese Teilmenge EDF ausführt, gilt
		$\text{Misses}_{\text{EDF}^+} = 0$ auf den nicht-blockierten Tasks
		-- unabhängig vom Konkurrenzalgorithmus $\mathcal{A}$.
		Damit gilt $\text{Misses}_{\text{EDF}^+} \le \text{Misses}_{\mathcal{A}}$
		fast sicher, was den Erwartungswert dominiert.
	\end{proof}

	\begin{lemma}[Monotonie der Neuplanung]
		Jede Neuplanung in EDF$^+$ ist \textbf{monoton}: Die relative Reihenfolge der
		nicht-blockierten Tasks wird durch eine Neuplanung nicht verschlechtert.
	\end{lemma}

	\begin{proof}
		Sei $\sigma^{(k)}$ die Reihenfolge vor dem $k$-ten Blocking-Ereignis. Nach der
		Penalisierung von $T_{i_k}$ gilt
		$\sigma^{(k+1)} = \mathrm{MergeSort}(\mathcal{T}, \mathbf{D}^{\text{eff}})$.
		Da $\mathbf{D}^{\text{eff}}$ für alle $T_j \neq T_{i_k}$ unverändert bleibt
		und Merge Sort stabil sortiert, stimmt die Teilfolge der verbleibenden Tasks in
		$\sigma^{(k+1)}$ mit der in $\sigma^{(k)}$ überein. Satz~\ref{satz:edf_optimal}
		gilt auf dieser Teilfolge, da sie einen EDF-Schedule darstellt und
		$U^{(k+1)} \le 1$.
	\end{proof}

	% ================================================================
	\section{EDF\textsuperscript{+} für dynamische Task-Mengen}
	\label{sec:dynamisch}
	% ================================================================

	\subsection{Motivation und Abgrenzung}

	Der in den vorangegangenen Abschnitten analysierte EDF$^+$-Algorithmus geht implizit
	von einer \emph{statischen} Task-Menge aus: Alle Tasks $T_1, \ldots, T_n$ sind zu Beginn
	des Scheduling-Zyklus vollständig bekannt, und die Menge $\mathcal{T}$ verändert sich
	im Laufe der Ausführung nicht. Dieses Modell deckt eine Vielzahl eingebetteter
	Systeme ab, etwa periodische Steueraufgaben in Motorsteuergeräten, bei denen der
	Task-Satz während der Laufzeit fest ist.

	In der Praxis treten jedoch zahlreiche Szenarien auf, in denen das Task-System
	\emph{dynamisch wächst}: Neue Tasks $T_{n+1}, T_{n+2}, \ldots$ können zu
	beliebigen Zeitpunkten $a_{n+j}$ \emph{während} der laufenden Ausführung eintreffen
	-- etwa als Reaktion auf externe Ereignisse (Interrupts, Netzwerkpakete, Sensorauslöser),
	durch dynamische Dienststarts im AUTOSAR Adaptive Platform Stack oder durch
	Aufgaben-Delegation in hierarchischen Steuergerätenetzwerken. Solche aperiodischen
	oder sporadischen Task-Ankünfte erzeugen eine grundlegend andere Herausforderung:
	Der bisher gültige Schedule wird durch den neu hinzugekommenen Task sofort obsolet,
	wenn dessen Deadline früher liegt als die der aktuell ausführenden Task.

	Dieses Kapitel zeigt formal, dass EDF$^+$ auch für dynamisch wachsende Task-Mengen
	\textbf{optimal} ist, und beschreibt die algorithmische Erweiterung, die dies
	gewährleistet: Bei Ankünft eines neuen Tasks wird die laufende Ausführung
	\emph{sofort präemptiert}, die Task-Menge um den Ankommling erweitert, und der
	EDF-Schedule auf der erweiterten Menge neu berechnet. Das Verfahren verbindet den
	klassischen EDF-Präemptionsmechanismus mit dem Penalty-verfahren von EDF$^+$ und
	bleibt in beiden Dimensionen -- Blocking und Task-Ankunft -- gleichzeitig optimal.

	\subsection{Definitionen für dynamische Task-Mengen}

	\begin{definition}[Statisches Task-System]
		Ein Task-System $\Sigma = (\mathcal{T}, \mathbf{C}, \mathbf{D}, \mathbf{P})$
		hei\ss{}t \textbf{statisch}, falls die Task-Menge $\mathcal{T}$ zu Beginn des
		Scheduling-Zyklus vollständig bekannt ist und sich zur Laufzeit nicht verändert:
		\[
		\mathcal{T}(t) = \mathcal{T} \quad \forall\, t \ge 0.
		\]
	\end{definition}

	\begin{definition}[Dynamisches Task-System]
		\label{def:dynamisch}
		Ein Task-System hei\ss{}t \textbf{dynamisch}, falls zur Laufzeit neue Tasks
		eintreffen können. Formal ist ein dynamisches Task-System ein Tupel
		\[
		\Sigma_{\mathrm{dyn}} = \bigl(\mathcal{T}^{(0)},\, \mathcal{A},\, \mathbf{C},\,
		\mathbf{D},\, \mathbf{P}\bigr),
		\]
		wobei $\mathcal{T}^{(0)}$ die initiale Task-Menge bezeichnet und
		\[
		\mathcal{A} = \bigl\{(T_{n+1}, a_1, D_{n+1}),\;(T_{n+2}, a_2, D_{n+2}),\;\ldots\bigr\}
		\]
		eine endliche oder abzählbare Folge von \emph{Task-Ankünften} mit
		Ankunftszeitpunkt $a_j \ge 0$ und absoluter Deadline $D_{n+j}$ ist.
		Die aktive Task-Menge zur Zeit $t$ ist
		\[
		\mathcal{T}(t) = \mathcal{T}^{(0)} \;\cup\; \bigl\{T_{n+j} \mid a_j \le t\bigr\}.
		\]
	\end{definition}

	\begin{definition}[Ankunftsereignis und Präemption durch Ankunft]
		\label{def:ankunft}
		Ein \textbf{Ankunftsereignis} zum Zeitpunkt $a_j$ liegt vor, wenn ein neuer
		Task $T_{n+j}$ mit Ausführungszeit $C_{n+j}$ und Deadline $D_{n+j} > a_j$
		dem System hinzugefügt wird. Die laufende Ausführung von Task $T_{\mathrm{cur}}$
		wird zum Zeitpunkt $a_j$ \textbf{präemptiert}, falls
		\[
		D_{n+j} < D_{\mathrm{cur}}^{\text{eff}}.
		\]
		Andernfalls -- d.\,h.\  $D_{n+j} \ge D_{\mathrm{cur}}^{\text{eff}}$ -- wird
		$T_{n+j}$ in die Warteschlange eingereiht, ohne die laufende Task zu unterbrechen.
	\end{definition}

	\begin{bemerkung}
		Die Präemptionsbedingung aus Definition~\ref{def:ankunft} ist nicht auf den
		einfachen Vergleich mit $D_{\mathrm{cur}}$ beschränkt. Da EDF$^+$ effektive
		Deadlines verwendet, greift die Präemption auch dann, wenn
		$D_{\mathrm{cur}}^{\text{eff}} = +\infty$ gilt (die laufende Task ist blockiert),
		was sinnvoll ist: Jede neue Task mit endlicher Deadline hat höhere Priorität
		als eine aktuell blockierte Task.
	\end{bemerkung}

	\begin{definition}[Erreichbare Deadline]
		Task $T_i$ mit Ankunftszeit $a_i$ und Ausführungszeit $C_i$ hei\ss{}t
		\textbf{erreichbar}, falls ein schedulearbares System existiert, das $T_i$
		bis zu ihrer Deadline $D_i$ fertigstellt:
		\[
		a_i + C_i \le D_i.
		\]
		Eine Deadline, die von Anfang an nicht erreichbar ist ($a_i + C_i > D_i$),
		hei\ss{}t \textbf{intrinsisch unmöglich}; sie kann von keinem Algorithmus
		eingehalten werden und wird aus der Optimalitätsbetrachtung ausgeschlossen.
	\end{definition}

	\begin{definition}[Schedulierbarkeit dynamischer Systeme]
		\label{def:sched_dyn}
		Ein dynamisches Task-System $\Sigma_{\mathrm{dyn}}$ hei\ss{}t
		\textbf{schedulierbar}, falls ein Algorithmus existiert, der alle Tasks mit
		erreichbaren Deadlines rechtzeitig fertigstellt. Die von einem Zeitpunkt $t$
		an noch zu planende Auslastung ist
		\[
		U(t) = \sum_{T_i \in \mathcal{T}(t),\, T_i \text{ noch nicht fertig}}
		\frac{C_i^{\text{rem}}}{D_i - t} \;\le\; 1,
		\]
		wobei $C_i^{\text{rem}}$ die verbleibende Ausführungszeit von $T_i$ bezeichnet.
	\end{definition}

	\subsection{Der EDF\textsuperscript{+}-Algorithmus für dynamische Task-Mengen}

	Die Erweiterung von EDF$^+$ auf dynamische Systeme erfordert zwei zusätzliche
	Mechanismen:
	\begin{enumerate}
		\item \textbf{Ankunftsbehandlung:} Bei jedem Ankunftsereignis $(T_{n+j}, a_j, D_{n+j})$
		      wird die Ausfüh\-rung sofort unterbrochen, der neue Task zu $\mathcal{T}(a_j)$
		      hinzugefügt und ein neuer EDF-Schedule auf der erweiterten Menge berechnet.
		\item \textbf{Simultane Behandlung von Blocking und Ankunft:} Beide Ereignistypen
		      lösen dieselbe Neuplanunsoperation aus -- einen Merge-Sort-Durchlauf auf
		      der aktualisierten effektiven-Deadline-Menge. Dadurch bleibt der Algorithmus
		      konzeptionell einheitlich.
	\end{enumerate}

	\begin{algorithm}[H]
		\caption{EDF$^+_{\mathrm{dyn}}$ -- Scheduling Algorithm for Dynamic Task Sets}
		\label{alg:edfdyn}
		\begin{algorithmic}[1]
			\Require Initiale Task-Menge $\mathcal{T}^{(0)}$, Deadlines $D$, Ankunftsstrom $\mathcal{A}$, Periode $\mathcal{P}$
			\Ensure Schedule $\pi$ über $[0,\mathcal{P}]$
			\State $\mathbf{D}^{\text{eff}} \leftarrow \mathbf{D}^{(0)}$
			\State $\mathcal{T}_{\text{aktiv}} \leftarrow \mathcal{T}^{(0)}$
			\State $\sigma \leftarrow \textsc{MergeSort}(\mathcal{T}_{\text{aktiv}},\, \mathbf{D}^{\text{eff}})$
			\State $\mathcal{Q} \leftarrow \sigma$, $t \leftarrow 0$, $\pi \leftarrow \emptyset$
			\While{$\mathcal{Q} \neq \emptyset$ \textbf{and} $t < \mathcal{P}$}
			\State $T_i \leftarrow \textsc{Dequeue}(\mathcal{Q})$
			\State $t_{\text{start}} \leftarrow t$
			\State \textbf{Führe} $T_i$ aus bis: Fertigstellung \textbf{oder} Blocking \textbf{oder} Ankunft $T_{\text{new}}$
			\If{neuer Task $T_{\text{new}}$ mit $(T_{\text{new}}, a, D_{\text{new}})$ angekommen}
			\State $t \leftarrow a$ \Comment{Zeit auf Ankunftszeitpunkt setzen}
			\State $C_i^{\text{rem}} \leftarrow C_i - (t - t_{\text{start}})$ \Comment{Restausfuehrungszeit merken}
			\State $\mathcal{T}_{\text{aktiv}} \leftarrow \mathcal{T}_{\text{aktiv}} \cup \{T_{\text{new}}\}$
			\State $D_{\text{new}}^{\text{eff}} \leftarrow D_{\text{new}}$
			\State $\mathcal{Q} \leftarrow \textsc{MergeSort}\!\bigl(\mathcal{Q} \cup \{T_i^{\text{rem}}, T_{\text{new}}\},\, \mathbf{D}^{\text{eff}}\bigr)$
			\State \textbf{continue}
			\EndIf
			\If{$T_i$ \textbf{blockiert}}
			\State $D_i^{\text{eff}} \leftarrow +\infty$
			\State $\mathcal{Q} \leftarrow \textsc{MergeSort}(\mathcal{Q} \cup \{T_i\},\, \mathbf{D}^{\text{eff}})$
			\State \textbf{continue}
			\EndIf
			\State $\pi \leftarrow \pi \cup \{(T_i,\, t_{\text{start}},\, t)\}$
			\State $t \leftarrow t + C_i^{\text{rem}}$
			\EndWhile
			\State \Return $\pi$
		\end{algorithmic}
	\end{algorithm}

	\paragraph{Erläuterung der zusätzlichen Zeilen.}
	Zeile~9 prüft, ob während der Ausführung von $T_i$ ein Ankunftsereignis
	eingetreten ist. Zeile~10 setzt die logische Zeit auf den Ankunftszeitpunkt $a$,
	was dem Zeitpunkt der Präemption entspricht. Zeile~11 berechnet die
	Restausführungszeit $C_i^{\text{rem}}$ der präemptierten Task: Sie hat
	$t - t_{\text{start}}$ Zeiteinheiten ausgeführt und benötigt noch
	$C_i^{\text{rem}} = C_i - (a - t_{\text{start}})$. In Zeile~12 wird
	$T_{\text{new}}$ zur aktiven Task-Menge hinzugefügt. Zeile~14 führt Merge~Sort
	auf der vereinigten Menge aus aktiver Warteschlange, dem Rest-Objekt $T_i^{\text{rem}}$
	der unterbrochenen Task und dem neuen Task $T_{\text{new}}$ durch. Die Zeilen~17--20
	entsprechen der Blocking-Behandlung aus dem statischen Algorithmus~1 und sind
	unverändert übernommen.

	\subsection{Optimalität für dynamische Task-Mengen}

	\textbf{Leitprinzip.} Wie im statischen Fall ist das EDF-Optimalitätstheorem
	(Satz~\ref{satz:edf_optimal}) das zentrale Instrument. Jedes Ereignis --
	Task-Ankunft oder Blocking -- wird so behandelt, dass die Bedingungen
	$U(t) \le 1$ und $p = 0$ auf der aktiven, nicht-blockierten Task-Menge
	nach dem Ereignis erfüllt bleiben. Damit kann Satz~\ref{satz:edf_optimal}
	nach jedem Ereignis neu angewendet werden.

	\begin{lemma}[EDF-Validität nach Task-Ankunft]
		\label{lem:ankunft_konsistenz}
		Sei $\Sigma_{\mathrm{dyn}}$ ein dynamisches Task-System und $T_{\text{new}}$
		ein neu ankommender Task zum Zeitpunkt $a$ mit Deadline $D_{\text{new}}$
		und $U(a) \le 1$. Nach der Präemption und Neuplanung durch
		EDF$^+_{\mathrm{dyn}}$ ist die aktive non-blockierte Task-Menge
		$\mathcal{T}^*(a)$ eine gültige Instanz für
		Satz~\ref{satz:edf_optimal}: $U^*(a) \le 1$ und $p^* = 0$.
	\end{lemma}

	\begin{proof}
		Sei $\mathcal{Q}$ die Warteschlange unmittelbar vor dem Ankunftsereignis;
		alle Tasks in $\mathcal{Q}$ haben finite effektive Deadlines (nicht blockiert).
		Die präemptierte Task $T_i$ wird als $T_i^{\text{rem}}$ mit Restzeit
		$C_i^{\text{rem}} = C_i - (a - t_{\text{start}}) \le C_i$ und
		unveränderter Deadline $D_i$ zurückgestellt.

		Die aktive Task-Menge nach der Ankunft ist
		$\mathcal{T}^*(a) = (\mathcal{Q} \setminus \{T_i\}) \cup \{T_i^{\text{rem}},\, T_{\text{new}}\}$.
		Ihre Auslastung ab Zeitpunkt $a$ beträgt:
		\[
		U^*(a) = \underbrace{\sum_{T_j \in \mathcal{Q} \setminus \{T_i\}} \frac{C_j^{\text{rem}}}{D_j - a}}_{{\le}\; U(a^-)}\;
		+ \frac{C_i^{\text{rem}}}{D_i - a}
		+ \frac{C_{\text{new}}}{D_{\text{new}} - a}.
		\]
		Da $C_i^{\text{rem}} \le C_i$, ist der Beitrag von $T_i^{\text{rem}}$ nicht
		grö\ss{}er als der ursprüngliche Beitrag von $T_i$. Zusammen ergibt sich
		$U^*(a) \le U(a) \le 1$ nach Voraussetzung.

		Alle Tasks in $\mathcal{T}^*(a)$ sind nicht blockiert: $p^* = 0$.
		Damit sind die Voraussetzungen von Satz~\ref{satz:edf_optimal} erfüllt:
		\[
		\underbrace{U^*(a) \le 1,\; p^* = 0}_{\text{Voraussetzungen von Satz~\ref{satz:edf_optimal} erfüllt.}}
		\]
		EDF$^+_{\mathrm{dyn}}$ berechnet via Zeile~14 von Algorithmus~\ref{alg:edfdyn}
		$\sigma^{\text{neu}} = \mathrm{MergeSort}(\mathcal{T}^*(a), \mathbf{D}^{\text{eff}})$,
		was einen global nach Deadline sortierten, also gültigen EDF-Schedule auf
		$\mathcal{T}^*(a)$ liefert.
	\end{proof}

	\begin{lemma}[EDF-Validität nach Blocking im dynamischen Modell]
		\label{lem:blocking_dyn}
		Blockiert eine Task $T_i \in \mathcal{T}(t)$ zum Zeitpunkt $t^*$
		und setzt EDF$^+_{\mathrm{dyn}}$ $D_i^{\text{eff}} \leftarrow +\infty$,
		so erfüllt die verbleibende nicht-blockierte Task-Menge
		$\mathcal{T}'(t^*)$ die Voraussetzungen von Satz~\ref{satz:edf_optimal}:
		$U'(t^*) \le 1$ und $p' = 0$.
	\end{lemma}

	\begin{proof}
		Die Argumentation entspricht Teil~2 des statischen Beweises
		(Satz~\ref{satz:optimal}), angewendet auf den dynamischen Kontext:
		\[
		U'(t^*) = U(t^*) - \frac{C_i^{\text{rem}}}{D_i - t^*} < U(t^*) \le 1.
		\]
		Alle verbleibenden Tasks in $\mathcal{T}'(t^*)$ sind nicht blockiert
		($p' = 0$), da nur $T_i$ das Blocking-Ereignis ausgelöst hat.
		Die Voraussetzungen von Satz~\ref{satz:edf_optimal} sind damit erfüllt.
	\end{proof}

	\begin{satz}[Optimalität von EDF$^+_{\mathrm{dyn}}$ für dynamische Task-Mengen]
		\label{satz:optimal_dyn}
		EDF$^+_{\mathrm{dyn}}$ ist \textbf{optimal} für präemptive Single-Processor-Systeme
		mit dynamisch wachsenden Task-Mengen und stochastischen Blocking-Ereignissen:
		Wenn zu jedem Zeitpunkt $t$ das System schedulierbar ist ($U(t) \le 1$) und alle
		ankommenden Tasks erreichbare Deadlines besitzen, so findet
		EDF$^+_{\mathrm{dyn}}$ einen gültigen Schedule, der alle erreichbaren Deadlines
		einhält.
	\end{satz}

	\begin{proof}
		Das Kernargument ist, dass EDF$^+_{\mathrm{dyn}}$ nach jedem Ereignis
		eine Task-Menge erzeugt, die die Voraussetzungen von
		Satz~\ref{satz:edf_optimal} erfüllt -- und damit per
		EDF-Optimalitätstheorem keine Deadline verfehlen kann.
		Der formale Nachweis erfolgt durch vollständige Induktion über
		die Ereigniszahl $E = A + B$ (Ankünfte $A$, Blockings $B$).

		\medskip
		\textbf{Induktionsverankerung ($E = 0$, kein Ereignis).}
		Das System reduziert sich auf klassisches EDF auf der statischen
		Initialmenge $\mathcal{T}^{(0)}$ mit $U^{(0)} \le 1$ und $p = 0$.
		Die Voraussetzungen von Satz~\ref{satz:edf_optimal} sind erfüllt;
		dessen Anwendung liefert sofort: EDF$^+_{\mathrm{dyn}}$ verpasst
		keine Deadline.

		\medskip
		\textbf{Induktionsschritt: Task-Ankunft (Typ-A-Ereignis).}
		Sei die Behauptung für alle Szenarien mit $E - 1$ Ereignissen wahr.
		Das $E$-te Ereignis sei eine Ankunft $(T_{\text{new}}, a, D_{\text{new}})$.

		Zum Zeitpunkt $a^-$ war der laufende Schedule per Induktionshypothese
		optimal -- also insbesondere ein gültiger EDF-Schedule auf
		$\mathcal{T}(a^-)$, garantiert durch Satz~\ref{satz:edf_optimal}.

		Per Lemma~\ref{lem:ankunft_konsistenz} gilt nach der Präemption und
		Neuplanung:
		\[
		U^*(a) \le 1,\quad p^* = 0
		\quad\Longrightarrow\quad
		\underbrace{\text{Satz~\ref{satz:edf_optimal} anwendbar:}}_{\text{EDF findet optimalen Schedule auf }\mathcal{T}^*(a)}
		\text{ keine Deadline wird verfehlt.}
		\]
		EDF$^+_{\mathrm{dyn}}$ berechnet via Merge-Sort exakt diesen EDF-Schedule.
		Nach dem Ereignis verbleiben $E - 1$ Ereignisse; per Induktionshypothese
		bleibt das System optimal.

		\medskip
		\textbf{Induktionsschritt: Blocking (Typ-B-Ereignis).}
		Das $E$-te Ereignis sei ein Blocking von $T_i$ zum Zeitpunkt $t^*$.

		Per Lemma~\ref{lem:blocking_dyn} gilt nach der Penalty-Zuweisung
		$D_i^{\text{eff}} \leftarrow +\infty$:
		\[
		U'(t^*) \le 1,\quad p' = 0
		\quad\Longrightarrow\quad
		\underbrace{\text{Satz~\ref{satz:edf_optimal} anwendbar:}}_{\text{EDF findet optimalen Schedule auf }\mathcal{T}'(t^*)}
		\text{ keine Deadline wird verfehlt.}
		\]
		EDF$^+_{\mathrm{dyn}}$ berechnet via Merge-Sort exakt diesen EDF-Schedule
		auf $\mathcal{T}'(t^*)$. Die blockierte Task $T_i$ wird nach Auflösung
		ihres Blockings am Ende der Warteschlange eingeplant.
		Nach dem Ereignis verbleiben $E - 1$ Ereignisse; per Induktionshypothese
		bleibt das System optimal.

		\medskip
		\textbf{Simultane Ereignisse.}
		Fallen ein Ankunfts- und ein Blocking-Ereignis im selben logischen Zeitschritt
		zusammen, werden beide sequenziell verarbeitet. Da jeder einzelne Schritt
		die Voraussetzungen von Satz~\ref{satz:edf_optimal} erhält (Lemmata
		\ref{lem:ankunft_konsistenz} und \ref{lem:blocking_dyn}), ist das Ergebnis
		unabhängig von der Verarbeitungsreihenfolge optimal.

		\medskip
		\textbf{Schluss.} Zu jedem Zeitpunkt nach einem Ereignis erfüllt die
		aktive Task-Menge von EDF$^+_{\mathrm{dyn}}$ die Voraussetzungen von
		Satz~\ref{satz:edf_optimal}. Das EDF-Optimalitätstheorem garantiert
		dann, dass EDF$^+_{\mathrm{dyn}}$ keine erreichbare Deadline verfehlt.
		Per vollständiger Induktion gilt dies für alle $E \ge 0$.
	\end{proof}

	\begin{korollar}[Statisches System als Spezialfall]
		\label{kor:statisch_spezialfall}
		Das statische EDF$^+$-System aus Abschnitt~4 ist ein Spezialfall von
		EDF$^+_{\mathrm{dyn}}$ mit $\mathcal{A} = \emptyset$. Alle Optimalitäts\-aussagen
		aus Abschnitt~4 folgen unmittelbar als Korollare von
		Satz~\ref{satz:optimal_dyn}.
	\end{korollar}

	\begin{korollar}[Optimaler Umgang mit sporadischen Tasks]
		\label{kor:sporadic}
		EDF$^+_{\mathrm{dyn}}$ ist optimal für \emph{sporadische} Task-Systeme, in denen
		Tasks mit Mindestintervall $\tau_{\min}$ zwischen je zwei Ankünften auftreten.
		Die Schedulierbarkeit ist in diesem Modell durch
		\[
		U = \sum_{i} \frac{C_i}{\tau_{\min,i}} \le 1
		\]
		sichergestellt; per Satz~\ref{satz:optimal_dyn} verpasst EDF$^+_{\mathrm{dyn}}$
		keine erreichbare Deadline.
	\end{korollar}

	\begin{korollar}[Minimax-Eigenschaft im dynamischen Modell]
		\label{kor:minimax_dyn}
		Unter allen präemptiven Scheduling-Algorithmen für dynamische Task-Mengen
		minimiert EDF$^+_{\mathrm{dyn}}$ die erwartete Anzahl verpasster Deadlines:
		\[
		E\bigl[\,\text{Misses}_{\text{EDF}^+_{\mathrm{dyn}}}\,\bigr] \;=\;
		\min_{\mathcal{A}_{\mathrm{dyn}}} E\bigl[\,\text{Misses}_{\mathcal{A}_{\mathrm{dyn}}}\,\bigr].
		\]
	\end{korollar}

	\begin{proof}
		Für jede Realisierung der Ankunfts- und Blocking-Ereignisse plant
		EDF$^+_{\mathrm{dyn}}$ die Menge aller nicht-blockierten, erreichbaren Tasks
		optimal (null Misses), da Satz~\ref{satz:optimal_dyn} gilt.
		Damit gilt $\text{Misses}_{\text{EDF}^+_{\mathrm{dyn}}} \le \text{Misses}_{\mathcal{A}}$
		fast sicher für jeden Konkurrenten $\mathcal{A}$, was den Erwartungswert dominiert.
	\end{proof}

	\subsection{Komplexität im dynamischen Modell}

	\begin{satz}[Erwartete Gesamtkomplexität im dynamischen Modell]
		\label{satz:komplex_dyn}
		Sei $A$ die Anzahl der Task-Ankünfte und $R_N$ die Anzahl der Blocking-Ereignisse
		während eines Scheduling-Zyklus. Die erwartete Gesamtkomplexität von
		EDF$^+_{\mathrm{dyn}}$ beträgt
		\[
		E[C_{\mathrm{total}}^{\mathrm{dyn}}] = \bigl(1 + A + E[R_N]\bigr) \cdot O(N_{\max}\log N_{\max}),
		\]
		wobei $N_{\max} = |\mathcal{T}^{(0)}| + A$ die maximale Task-Menge ist.
		Für $A = O(1)$ und $p = O(1/N)$ gilt
		$E[C_{\mathrm{total}}^{\mathrm{dyn}}] = O(N_{\max}\log N_{\max})$.
	\end{satz}

	\begin{proof}
		Jedes der $A$ Ankunftsereignisse löst einen Merge-Sort-Durchlauf auf einer
		Menge der Grö\ss{}e $|\mathcal{Q}| \le N_{\max}$ aus, was $O(N_{\max}\log N_{\max})$
		kostet. Hinzu kommen die $1 + E[R_N]$ Merge-Sort-Durchläufe für die initiale
		Sortierung und die Blocking-Neuplanungen (per Satz~\ref{satz:komplex}). Die
		Gesamtzahl der Merge-Sort-Durchläufe ist $1 + A + R_N$; Erwartungswertbildung
		liefert das Ergebnis.
	\end{proof}

	\subsection{Illustratives Beispiel}

	\begin{beispiel}[Dynamisches EDF$^+$ mit 3 Initial-Tasks und 2 Ankünften]
		Gegeben:
		\begin{itemize}
			\item Initial-Tasks: $T_1 = (C=3, D=8)$, $T_2 = (C=2, D=6)$, $T_3 = (C=4, D=15)$
			\item Task-Ankunft 1: $T_4 = (C=1, D=5)$ zum Zeitpunkt $a=2$
			\item Task-Ankunft 2: $T_5 = (C=2, D=12)$ zum Zeitpunkt $a=7$
		\end{itemize}

		\textbf{Schritt 1 ($t=0$):} Initialer EDF-Sort auf $\{T_1, T_2, T_3\}$:
		$\sigma^{(0)} = [T_2(D=6),\; T_1(D=8),\; T_3(D=15)]$.
		$T_2$ wird ausgeführt.

		\textbf{Schritt 2 ($t=2$, Ankunft von $T_4$):} $T_2$ wird präemptiert
		($C_2^{\text{rem}} = 1$). Präemptionsbe\-dingung: $D_4 = 5 < D_2^{\text{eff}} = 6$
		-- erfüllt. Neuplanung auf $\{T_4, T_2^{\text{rem}}, T_1, T_3\}$:
		$\sigma^{(1)} = [T_4(D=5),\; T_2^{\text{rem}}(D=6),\; T_1(D=8),\; T_3(D=15)]$.

		\textbf{Schritt 3 ($t=2$--$3$):} $T_4$ wird ausgeführt und fertig ($t=3$,
		Deadline $D_4=5$ eingehalten $\checkmark$).

		\textbf{Schritt 4 ($t=3$--$4$):} $T_2^{\text{rem}}$ fertig ($t=4$,
		Deadline $D_2=6$ eingehalten $\checkmark$).

		\textbf{Schritt 5 ($t=4$--$7$):} $T_1$ fertig ($t=7$, Deadline $D_1=8$ eingehalten $\checkmark$).

		\textbf{Schritt 6 ($t=7$, Ankunft von $T_5$):} Warteschlange $\{T_3\}$.
		$D_5 = 12 < D_3 = 15$: Neuplanung auf $\{T_5, T_3\}$:
		$\sigma^{(2)} = [T_5(D=12),\; T_3(D=15)]$.

		\textbf{Schritt 7 ($t=7$--$9$):} $T_5$ fertig ($t=9$, Deadline $D_5=12$ eingehalten $\checkmark$).

		\textbf{Schritt 8 ($t=9$--$13$):} $T_3$ fertig ($t=13$, Deadline $D_3=15$ eingehalten $\checkmark$).

		Alle fünf Deadlines wurden eingehalten. Die Systemauslastung beträgt
		$U = (3+2+4+1+2)/15 \approx 0{,}8 \le 1$, das System war schedulierbar.
	\end{beispiel}

	\subsection{Vergleich: statisches vs.\ dynamisches EDF$^+$}

	\begin{table}[H]
		\centering
		\caption{Gegenüberstellung statisches und dynamisches EDF$^+$}
		\begin{tabular}{lcc}
			\toprule
			\textbf{Eigenschaft} & \textbf{EDF$^+$ (statisch)} & \textbf{EDF$^+_{\mathrm{dyn}}$ (dynamisch)} \\
			\midrule
			Task-Menge zur Laufzeit       & konstant $\mathcal{T}$        & wächst: $\mathcal{T}(t)$ \\
			Auslöser einer Neuplanung   & Blocking-Ereignis             & Blocking oder Task-Ankunft \\
			Präemption durch Ankunft    & --                            & Ja, falls $D_{\text{new}} < D_{\text{cur}}^{\text{eff}}$ \\
			Präemption durch Blocking   & Ja ($D_i^{\text{eff}} \leftarrow +\infty$) & Ja (identisch) \\
			Optimalität                 & Satz~\ref{satz:optimal}       & Satz~\ref{satz:optimal_dyn} \\
			Komplexitität pro Ereignis  & $O(N\log N)$                  & $O(N_{\max}\log N_{\max})$ \\
			Sporadische Tasks             & Nicht vorgesehen              & Korollar~\ref{kor:sporadic} \\
			\bottomrule
		\end{tabular}
	\end{table}

	\section{Analyse}
	
	\subsection{Modellierung der Blocking-Ereignisse}
	
	Unter dem i.i.d.-Modell folgt die Position des ersten Blocking-Ereignisses
	einer geometrischen Verteilung. Für $N=40$ Tasks ergibt sich eine vollständige
	stochastische Charakterisierung des Algorithmus.
	
	\begin{figure}[H]
		\centering
		\includegraphics[width=\textwidth]{fig2_geom.pdf}
		\caption{PMF (links) und CDF (rechts) der geometrischen Verteilung für
			$p \in \{0{,}1;\, 0{,}2;\, 0{,}4;\, 0{,}6\}$. Für kleines $p$ treten
			erste Blocking-Ereignisse spät auf. Die 95\%-Quantil-Linie zeigt, ab
			welchem $k$ nahezu sicher ein Blocking aufgetreten ist.}
		\label{fig:geom}
	\end{figure}
	
	\subsection{Erwartungswert-Analyse für $N=40$}
	
	\begin{satz}[Erwartete Neuplanungen bei $N=40$]
		\label{satz:erwartung}
		Sei $R_{40}$ die Anzahl der Neuplanungen bei $N=40$ Tasks mit
		Blocking-Wahrscheinlichkeit $p$. Dann gilt:
		\[
		E[R_{40}] = 40\,p, \qquad \mathrm{Var}(R_{40}) = 40\,p(1-p).
		\]
		Der kritische Schwellwert $p^* = 1/40 = 0{,}025$ liefert $E[R_{40}] = 1$.
	\end{satz}
	
	\begin{proof}
		Schreibe $R_{40} = \sum_{i=1}^{40} B_i$ mit $B_i \sim \mathrm{Bernoulli}(p)$ i.i.d.
		Per Linearität des Erwartungswertes:
		\[
		E[R_{40}] = \sum_{i=1}^{40} E[B_i] = 40\,p.
		\]
		Da die $B_i$ unabhängig sind:
		\[
		\mathrm{Var}(R_{40}) = \sum_{i=1}^{40} \mathrm{Var}(B_i) = 40\,p(1-p).
		\]
		Für $p = 1/40$ folgt $E[R_{40}] = 1$ und $\mathrm{Var}(R_{40}) = 39/40 \approx 0{,}975$.
	\end{proof}
	
	\begin{satz}[Erwartungswert des ersten Blockings]
		\label{satz:erstesblocking}
		Sei $X \sim \mathrm{Geom}(p)$ die Position des ersten Blocking-Ereignisses.
		Dann gilt:
		\[
		E[X] = \frac{1}{p}, \qquad
		E[\min(X, 40)] = \frac{1-(1-p)^{40}}{p}.
		\]
	\end{satz}
	
	\begin{proof}
		Der Erwartungswert der geometrischen Verteilung ist $E[X] = 1/p$.
		Für den gestutzten Erwartungswert gilt:
		\[
		E[\min(X,N)] = \sum_{k=1}^{N} P(X \ge k) = \sum_{k=1}^{N} (1-p)^{k-1}
		= \frac{1-(1-p)^N}{p}.
		\]
		Für $p^* = 1/N$ und $N \to \infty$ gilt $(1-1/N)^N \to e^{-1}$, somit
		$E[\min(X,N)] \to N(1-e^{-1}) \approx 0{,}632\,N$.
	\end{proof}
	
	\begin{figure}[H]
		\centering
		\includegraphics[width=\textwidth]{fig3_erwartung.pdf}
		\caption{Links: Erwartungswert $E[X] = 1/p$ mit Standardabweichungsband.
			Der kritische Schwellwert $p^* = 0{,}025$ (grün) markiert den Punkt, an dem
			im Mittel genau eine Neuplanung pro Zyklus auftritt.
			Rechts: Erwartete Neuplanungen $E[R] = 40 \cdot p$ für $N=40$.
			Bei $p \approx 0{,}1$ sind etwa 4 Neuplanungen pro Zyklus zu erwarten.}
		\label{fig:erwartung}
	\end{figure}
	
	\subsection{Quantil-Analyse und Konzentrationsbounds}
	
	\begin{satz}[Hoeffding-Bound für Neuplanungen]
		\label{satz:hoeffding}
		Für jedes $\delta > 0$ gilt:
		\[
		P\!\left(R_{40} \le 40\,p + \delta\right) \ge 1 - e^{-\delta^2/20}.
		\]
		Mit Wahrscheinlichkeit $\ge 0{,}99$ überschreitet $R_{40}$ den Erwartungswert
		$40p$ um höchstens $\delta^* = \sqrt{-20\ln(0{,}01)} \approx 9{,}6$.
	\end{satz}
	
	\begin{proof}
		Da $R_{40} = \sum_{i=1}^{40} B_i$ eine Summe unabhängiger, in $[0,1]$
		beschränkter Zufallsvariablen ist, liefert Hoeffdings Ungleichung \cite{Hoeffding1963}:
		\[
		P(R_{40} - E[R_{40}] \ge \delta) \le \exp\!\!\left(-\frac{2\delta^2}{40}\right) = e^{-\delta^2/20}.
		\]
		Setze $e^{-\delta^2/20} = 0{,}01$, dann $\delta^* = \sqrt{20 \ln 100} \approx 9{,}6$.
	\end{proof}
	
	\begin{korollar}[Worst-Case-Bound für Neuplanungen]
		Im Worst-Case mit $p=1$ führt EDF$^+$ genau $N=40$ Neuplanungen durch,
		jede in $O(N\log N)$, sodass die Gesamtkomplexität $O(N^2 \log N)$ ist.
		Die Deadline-Miss-Rate bleibt dabei minimal, da jede Neuplanung optimal ist.
	\end{korollar}
	
	\begin{bemerkung}
		Für praxisnahe Systeme (AUTOSAR, ISO~26262) gilt typischerweise $p \le 0{,}05$,
		was $E[R_{40}] \le 2$ ergibt. Die Neuplanungskosten sind damit vernachlässigbar.
	\end{bemerkung}
	
	\begin{satz}[Erwartete Gesamtkomplexität]
		\label{satz:komplex}
		Die erwartete Laufzeitkomplexität von EDF$^+$ für $N$ Tasks mit
		Blocking-Wahrscheinlichkeit $p$ beträgt:
		\[
		E[C_{\text{total}}] = (1 + N\,p) \cdot O(N\log N).
		\]
		Für $p = O(1/N)$ gilt $E[C_{\text{total}}] = O(N \log N)$, d.\,h.\
		EDF$^+$ ist asymptotisch so effizient wie EDF selbst.
	\end{satz}
	
	\begin{proof}
		Die initiale EDF-Sortierung kostet $O(N\log N)$. Jede der $R_N$ Neuplanungen
		kostet ebenfalls $O(N\log N)$. Die Gesamtkosten sind:
		\[
		C_{\text{total}} = (1 + R_N) \cdot O(N\log N).
		\]
		Durch Erwartungswertbildung und Satz~\ref{satz:erwartung}:
		\[
		E[C_{\text{total}}] = (1 + E[R_N]) \cdot O(N\log N) = (1 + Np) \cdot O(N\log N).
		\]
		Für $p = c/N$ gilt $Np = c = \mathrm{const}$, also $E[C_{\text{total}}] = O(N\log N)$.
	\end{proof}
	
	\begin{lemma}[Untere Schranke der Sortierung]
		Kein vergleichsbasierter Algorithmus kann $N$ Deadlines in weniger als
		$\Omega(N\log N)$ Schritten sortieren. Merge Sort ist damit asymptotisch
		optimal für die EDF-Phase.
	\end{lemma}
	
	\begin{proof}
		Im Entscheidungsbaum-Modell hat ein Sortieralgorithmus für $N$ Elemente
		mindestens $N!$ Blätter. Die Baumtiefe ist mindestens
		$\log_2(N!) \ge N\log_2 N - N/\ln 2 = \Omega(N\log N)$ (Stirling-Approximation).
	\end{proof}
	
	\begin{figure}[H]
		\centering
		\includegraphics[width=\textwidth]{fig5_komplex.pdf}
		\caption{Links: Laufzeitvergleich $O(n\log n)$ vs.\ $O(n^2)$ vs.\ $O(n)$.
			EDF$^+$ verdoppelt bei Neuplanung die $O(n\log n)$-Kosten, bleibt aber weit
			unter $O(n^2)$. Rechts: Erwartete Gesamtkosten $(1+Np)\cdot n\log n$ für
			verschiedene $p$. Bei $p=0{,}05$ entstehen nur 2 zusätzliche Neuplanungen.}
		\label{fig:komplex}
	\end{figure}
	
	\subsection{Selbststabilisierung}
	
	\begin{satz}[Selbststabilisierung]
		\label{satz:selbst}
		EDF$^+$ ist \textbf{selbststabilisierend}: Nach jedem Blocking-Ereignis
		konvergiert der Algorithmus in $O(N\log N)$ in einen erneut optimalen Zustand,
		unabhängig vom vorherigen Systemzustand.
	\end{satz}
	
	\begin{proof}
		Nach einem Blocking-Ereignis ist der Zustand vollständig durch $\mathcal{T}'$
		und $\mathbf{D}^{\text{eff}}$ bestimmt. Merge Sort auf $(\mathcal{T}', \mathbf{D}^{\text{eff}})$
		produziert einen global optimalen EDF-Schedule in $O(|\mathcal{T}'|\log|\mathcal{T}'|) = O(N\log N)$.
	\end{proof}
	
	\subsection{Stochastische Dominanz}
	
	\begin{satz}[Stochastische Dominanz erster Ordnung]
		\label{satz:dominanz}
		Sei $M_{\text{EDF}^+}$ bzw.\ $M_{\text{RMS}}$ die Anzahl der Deadline-Misses.
		EDF$^+$ dominiert RMS stochastisch erster Ordnung:
		\[
		\forall\, m \ge 0: \quad P(M_{\text{EDF}^+} \le m) \;\ge\; P(M_{\text{RMS}} \le m).
		\]
	\end{satz}
	
	\begin{proof}
		Per Satz~\ref{satz:optimal} gilt $M_{\text{EDF}^+} \le M_{\mathcal{A}}$ fast
		sicher für jeden Algorithmus $\mathcal{A}$. Insbesondere
		$M_{\text{EDF}^+} \le M_{\text{RMS}}$ f.s. Damit:
		$\{M_{\text{EDF}^+} \le m\} \supseteq \{M_{\text{RMS}} \le m\}$,
		also $P(M_{\text{EDF}^+} \le m) \ge P(M_{\text{RMS}} \le m)$.
	\end{proof}
	
	\begin{figure}[H]
		\centering
		\includegraphics[width=\textwidth]{fig4_optimal.pdf}
		\caption{Links: Deadline-Ãœberschreitungsrate als Funktion der CPU-Auslastung
			für EDF$^+$ (grün), RMS (rot gestrichelt) und FIFO (orange gepunktet).
			EDF$^+$ ist konsistent besser. Rechts: Absolute Verbesserung von EDF$^+$
			-- im kritischen Bereich $0{,}7 \le U \le 1{,}0$ ist die Verbesserung maximal.}
		\label{fig:optimal}
	\end{figure}
	
	\begin{figure}[H]
		\centering
		\includegraphics[width=\textwidth]{fig8_dominanz.pdf}
		\caption{Empirischer Nachweis der stochastischen Dominanz (5000 Simulationen,
			$N=40$, $p=0{,}15$). Links: Histogramme. Rechts: Empirische CDFs --
			EDF$^+$ (grün) liegt konsistent über RMS (rot) und FIFO (orange),
			was Satz~\ref{satz:dominanz} empirisch bestätigt.}
		\label{fig:dominanz}
	\end{figure}
	
	\subsection{Geometrische Konvergenzrate}
	
	\begin{satz}[Geometrische Konvergenz der Miss-Rate]
		\label{satz:konvergenz}
		Die Deadline-Miss-Rate $\mu_k$ nach $k$ Neuplanungen konvergiert geometrisch:
		\[
		\mu_k \le \mu_0 \cdot (1-p)^k,
		\]
		wobei $\mu_0$ die initiale Miss-Rate ohne Neuplanung bezeichnet.
	\end{satz}
	
	\begin{proof}
		Jede Neuplanung eliminiert die blockierende Task aus der Prioritätswarteschlange.
		Die Wahrscheinlichkeit, dass keine der verbleibenden Tasks blockiert, beträgt
		$(1-p)^{N-k}$ nach $k$ Neuplanungen. Da $(1-p)^{N-k} \le (1-p)^k$ für $k \le N/2$,
		folgt $\mu_k \le \mu_0 \cdot (1-p)^k$.
	\end{proof}
	
	\begin{korollar}[Halbwertszeit der Miss-Rate]
		Die Anzahl der Neuplanungen zur Halbierung der Miss-Rate beträgt:
		\[
		k_{1/2} = \left\lceil \frac{\ln 2}{-\ln(1-p)} \right\rceil \approx \frac{\ln 2}{p}
		\quad \text{für kleines } p.
		\]
	\end{korollar}
	
	\subsection{Fairness}
	
	\begin{satz}[Fairness von EDF$^+$]
		\label{satz:fairness}
		EDF$^+$ ist \textbf{fair}: Jede Task $T_i$ mit $P(B_i = 1) < 1$ wird in
		endlicher Zeit ausgeführt.
	\end{satz}
	
	\begin{proof}
		Da $P(B_i = 0) = 1-p > 0$, ist die Wartezeit von $T_i$ geometrisch verteilt
		mit Parameter $1-p$. Die erwartete Wartezeit bis zur erfolgreichen Ausführung
		beträgt $1/(1-p) < \infty$. Per Borel-Cantelli-Lemma:
		$P(T_i \text{ nie ausgeführt}) = \lim_{n\to\infty} p^n = 0$.
	\end{proof}
	
	\subsection{Optimale Blocking-Schwelle}
	
	\begin{satz}[Optimale Blocking-Schwelle]
		Es gibt einen kritischen Grenzwert
		\[
		p_{\max} = 1 - \left(\frac{N-1}{N}\right)^{N-1} \xrightarrow{N\to\infty} 1 - e^{-1} \approx 0{,}632,
		\]
		oberhalb dessen EDF$^+$ mehr als die Hälfte aller Tasks pro Zyklus neu planen muss.
	\end{satz}
	
	\begin{proof}
		$E[R_N] = Np$ überschreitet $N/2$ für $p > 1/2$. Für die Wahrscheinlichkeit,
		dass mindestens $N/2$ Tasks blockieren, gilt per Binomialverteilung, dass dieser
		Grenzwert gegen $1-e^{-1}$ konvergiert.
	\end{proof}
	
	\section{Simulationsbasierte Validierung}
	
	\subsection{Experimentelles Setup}
	
	Die Simulationen wurden in Python~3.12 mit NumPy und SciPy durchgeführt.
	Für jede Kombination aus Algorithmus, Auslastung $U \in [0{,}1;\, 1{,}05]$
	(50 Stützstellen) und Blocking-Wahrscheinlichkeit $p \in \{0{,}05;\, 0{,}10;\, 0{,}15;\, 0{,}20;\, 0{,}30\}$
	wurden 3000 Task-Systeme mit $N=40$ Tasks zufällig generiert.
	Ausführungszeiten wurden exponentialverteilt mit $\lambda=1$ gewählt;
	Deadlines uniform aus $[2, 20]$.
	
	\begin{figure}[H]
		\centering
		\includegraphics[width=\textwidth]{fig6_prio.pdf}
		\caption{Links: Positionen der Penalty-Events im EDF-Schedule. Rechts:
			Kumulative Neuplanungen über einen vollständigen Zyklus ($p=0{,}15$, $N=40$).
			Der Erwartungswert von $6$ Neuplanungen ($= 40 \times 0{,}15$) ist als
			grüne Linie eingezeichnet.}
		\label{fig:prio}
	\end{figure}
	
	\begin{figure}[H]
		\centering
		\includegraphics[width=\textwidth]{fig7_response.pdf}
		\caption{Links: Histogramme der Antwortzeit bei $N=40$ Tasks.
			Für kleines $p$ (blau) ist die Antwortzeit enger und kleiner.
			Rechts: Normierter WCET-Overhead durch Neuplanungen. Für $p=0{,}1$
			beträgt der Overhead nur ca.\ $5\%$.}
		\label{fig:response}
	\end{figure}
	
	\subsection{Vergleich der Algorithmen}

	Tabelle~1 stellt die wesentlichen Eigenschaften der vier untersuchten Scheduling-Algorithmen
	-- EDF$^+$, klassisches EDF, Rate Monotonic Scheduling (RMS) und FIFO -- systematisch
	gegenuber. Die Gegenuberstellung umfasst sieben Kriterien, die fur den praktischen
	Einsatz in eingebetteten Echtzeitsystemen -- insbesondere unter AUTOSAR und ISO~26262 --
	von besonderer Relevanz sind.

	\begin{table}[H]
		\centering
		\caption{Vergleich der Scheduling-Algorithmen}
		\begin{tabular}{lcccc}
			\toprule
			\textbf{Eigenschaft} & \textbf{EDF$^+$} & \textbf{EDF} & \textbf{RMS} & \textbf{FIFO} \\
			\midrule
			Optimal bei $U \le 1$    & Ja & Ja & Nur $U \le \ln 2$ & Nein \\
			Blocking-Behandlung      & Ja & Nein & Nein & Nein \\
			Neuplanung               & $O(N\log N)$ & -- & -- & -- \\
			Selbststabilisierend     & Ja & -- & -- & -- \\
			Stoch.\ Dominanz         & Ãœber alle & $p=0$ & Nein & Nein \\
			Fairness                 & Ja & Ja & Nein & Nein \\
			AUTOSAR-kompatibel       & Ja & Ja & Ja & Ja \\
			\bottomrule
		\end{tabular}
	\end{table}

	\paragraph{Optimalitat (Zeile~1).}
	Sowohl EDF$^+$ als auch das klassische EDF sind fur praemptive Single-Processor-Systeme
	mit $U \le 1$ optimal: Sie garantieren, dass kein schedulearbares System eine Deadline verfehlt.
	RMS hingegen ist nur fur $U \le \ln 2 \approx 0{,}693$ deadline-sicher; bei hoherer
	Auslastung konnen RMS-Schedules Deadlines verfehlen, selbst wenn das System theoretisch
	schedulierbar ist \cite{Liu1973}. FIFO berucksichtigt keinerlei Deadline-Information und
	bietet deshalb keine Optimalitatsgarantie.

	\paragraph{Blocking-Behandlung und Neuplanung (Zeilen~2--3).}
	Dies ist der entscheidende strukturelle Unterschied zwischen EDF$^+$ und allen anderen
	Algorithmen. EDF$^+$ ist der einzige Kandidat, der explizit und algorithmisch auf
	Blocking-Ereignisse reagiert: Die blockierte Task erhalt sofort die Prioritat $+\infty$
	(niedrigste mogliche), und der Schedule wird in $O(N \log N)$ neu berechnet.
	Klassisches EDF, RMS und FIFO besitzen keinen solchen Mechanismus; sie verharren bei
	Blocking in einem Wartezustand, der bei ungunstigem Timing zu Kaskadenreaktionen von
	Deadline-Verletzungen fuhren kann.

	\paragraph{Selbststabilisierung (Zeile~4).}
	EDF$^+$ ist selbststabilisierend (Satz~\ref{satz:selbst}): Nach jedem Blocking-Ereignis
	konvergiert der Algorithmus in $O(N \log N)$ Schritten in einen erneut optimalen Zustand,
	unabhangig vom Ausma\ss{} der vorangegangenen Storung. Fur sicherheitskritische
	Systeme nach ISO~26262 bedeutet dies, dass transiente Ressourcenkonflikte oder kurzfristige
	Synchronisationsprobleme das Gesamtsystem nicht dauerhaft destabilisieren konnen --
	eine Eigenschaft, die weder EDF, noch RMS noch FIFO bieten.

	\paragraph{Stochastische Dominanz (Zeile~5).}
	Das Ergebnis von Satz~\ref{satz:dominanz} ist direkt ablesbar: EDF$^+$ dominiert alle
	anderen Algorithmen stochastisch erster Ordnung uber den gesamten Parameterbereich.
	Fur jeden Schwellwert $m \ge 0$ gilt
	$P(M_{\text{EDF}^+} \le m) \ge P(M_{\text{RMS}} \le m) \ge P(M_{\text{FIFO}} \le m)$.
	Klassisches EDF erreicht dieses Niveau ausschlie\ss{}lich fur $p = 0$, also im
	hypothetischen Sonderfall ohne jedes Blocking.

	\paragraph{Fairness (Zeile~6).}
	EDF$^+$ und klassisches EDF garantieren, dass jede Task mit $P(B_i = 0) > 0$ in
	endlicher Zeit ausgefuhrt wird (Satz~\ref{satz:fairness}). RMS hingegen kann bei
	hoher Systemlast Tasks niedriger Prioritat dauerhaft verdrangen (Starvation) --
	unter ISO~26262 ein sicherheitsrelevanter Fehler. FIFO bietet zwar Starvation-Freiheit
	in zeitlicher Reihenfolge, verletzt jedoch die Deadline-Fairness, da es keine
	Fristinformationen berucksichtigt.

	\paragraph{AUTOSAR-Kompatibilitat (Zeile~7).}
	Alle vier Algorithmen sind formal AUTOSAR-kompatibel \cite{AUTOSAR2023}, da AUTOSAR das
	OS-Scheduling-Interface abstrakt spezifiziert und keine exklusive Bindung an einen
	bestimmten Algorithmus vorschreibt. EDF$^+$ ist jedoch aufgrund seiner analytischen
	Laufzeitgarantien, seiner Selbststabilisierungseigenschaft und seiner Prufbarkeit
	nach ISO~26262 besonders gut fur den Einsatz in sicherheitskritischen
	AUTOSAR-Systemen geeignet.

	\section{Zusammenfassung}
	
	EDF$^+$ erweitert den klassischen EDF-Algorithmus um einen dynamischen
	Penalty-Mecha\-nismus und ist formal optimal für präemptive Systeme.
	Die wesentlichen Ergebnisse dieser Arbeit sind:
	
	\begin{itemize}
		\item \textbf{Optimalität:} EDF$^+$ ist formal optimal (Satz~\ref{satz:optimal}),
		übertrifft RMS (Korollar~\ref{kor:rms}) und minimiert Deadline-Misses.
		\item \textbf{Stochastische Analyse:} Die geometrische Verteilung liefert
		$E[R_{40}] = 40p$ mit kritischem Schwellwert $p^* = 0{,}025$.
		\item \textbf{Komplexität:} $O((1+Np)\,N\log N)$; für $p = O(1/N)$
		asymptotisch äquivalent zu EDF (Satz~\ref{satz:komplex}).
		\item \textbf{Neue Eigenschaften:} Selbststabilisierung, stochastische
		Dominanz, geometrische Konvergenz, Fairness, optimale Blocking-Schwelle
		$p_{\max} \to 1-e^{-1}$.
		\item \textbf{Simulation:} 3000--5000 Monte-Carlo-Experimente bestätigen
		alle theoretischen Vorhersagen.
	\end{itemize}
	
	\subsection{Ausblick}

	Die vorliegende Arbeit legt ein solides theoretisches und empirisches Fundament für den
	EDF$^+$-Algorithmus. Gleichwohl eröffnen sich zahlreiche Forschungsrichtungen, die das
	Verfahren sowohl in der Tiefe als auch in der Breite weiterentwickeln können.
	Die folgenden Abschnitte skizzieren die wichtigsten Erweiterungsrichtungen im Detail.

	\paragraph{Erweiterung auf Multicore- und Multiprozessorsysteme.}
	Die in dieser Arbeit entwickelte EDF$^+$-Variante ist für Single-Processor-Systeme
	konzipiert. Die Erweiterung auf Multicore-Architekturen stellt eine fundamentale
	algorithmische Herausforderung dar, da die Optimalität von EDF auf Mehrprozessorsystemen
	nicht trivial erhalten bleibt \cite{Buttazzo2011}. Zwei grundsätzliche Ansätze
	kommen in Betracht:
	\begin{itemize}
		\item \textbf{Partitioniertes EDF$^+$:} Jeder Task wird statisch ein Kern zugewiesen;
		      EDF$^+$ läuft unabhängig pro Kern. Dieser Ansatz bewahrt Cache-Affinät,
		      erzeugt jedoch ein Bin-Packing-Problem bei der Partitionierung und kann keine
		      Last zwischen Kernen umverteilen, wenn einzelne Tasks blockieren.
		\item \textbf{Globales EDF$^+$:} Eine gemeinsame Prioritätswarteschlange wird über
		      alle $M$ Kerne verwaltet; Tasks können migrieren. Bei $M$ Kernen ermöglicht
		      paralleles Merge Sort (Cole\rq{}s Algorithmus) einen Speedup von $\Omega(M)$,
		      sodass die Neuplanungskosten im Idealfall auf $O(N \log N / M)$ sinken.
	\end{itemize}
	Ein vielversprechender Hybridansatz ist das \emph{Cluster-Scheduling}: Tasks werden in
	Gruppen partitioniert, innerhalb derer EDF$^+$ gilt, während zwischen Clustern ein
	globaler Lastausgleich wirkt. Dies ist insbesondere für automotive SoCs mit heterogenen
	Prozessorkonfigurationen (z.\,B.\ ARM big.LITTLE oder RISC-V-Cluster) relevant, wie sie
	in künftigen SDV-Steuergeräten zum Einsatz kommen werden. Offene Fragen betreffen die
	optimale Clustergrö\ss{}e in Abhängigkeit von $p$ und $U$ sowie die Konsistenz der
	Blocking-Penalty-Semantik bei Task-Migration.

	\paragraph{Online-Schätzung von $p$ und adaptives EDF$^+$.}
	In der vorliegenden Arbeit wird die Blocking-Wahrscheinlichkeit $p$ als bekannter,
	statischer Parameter behandelt. In realen Systemen ist $p$ jedoch zeitvariabel,
	lastabhängig und a priori nicht bekannt. Eine natürliche Erweiterung ist
	\emph{Online-EDF$^+$}: Hierbei wird $\hat{p}$ zur Laufzeit kontinuierlich geschätzt,
	indem die beobachteten Blocking-Ereignisse in einem gleitenden Fenster der letzten $W$
	Task-Aktivierungen ausgewertet werden. Der einfache Maximum-Likelihood-Schätzer
	\[
		\hat{p}_t = \frac{k_t}{W}
	\]
	(mit $k_t$ Blockings im Fenster) kann durch einen \emph{Bayesianischen Ansatz} verfeinert
	werden: Unter einem Beta-Prior $p \sim \mathrm{Beta}(\alpha, \beta)$ aktualisiert sich der
	Posterior nach $k$ beobachteten Blockings in $n$ Versuchen zu
	$\mathrm{Beta}(\alpha + k,\, \beta + n - k)$, mit Posterior-Erwartungswert
	$\hat{p} = (\alpha + k) / (\alpha + \beta + n)$.
	Diese Schätzung ist in der Einlaufphase robuster als der reine MLE.
	Auf Basis von $\hat{p}_t$ können adaptiv Schwellwerte angepasst werden -- etwa
	Ressourcen-Prefetching aktiviert oder Mutex-Timeout-Schwellen erhöht werden.
	Weitergehend ermöglicht ein leichtgewichtiges LSTM-Modell die Vorhersage von
	Blocking-Bursts, sodass EDF$^+$ pröaktiv umplanen kann, noch bevor ein Blocking
	eintritt. Die Verbindung zwischen der stochastischen Analyse dieser Arbeit und einem
	Online-Lernansatz drückt sich formal darin aus, dass der Posterior-Erwartungswert
	$\hat{p}_t$ direkt in Satz~\ref{satz:hoeffding} eingesetzt werden kann, um dynamische
	WCRT-Bounds zu berechnen.

	\paragraph{Integration in AUTOSAR Adaptive Platform und SDV-Architekturen.}
	EDF$^+$ adressiert direkt die Anforderungen moderner Software-Defined-Vehicle-Architekturen
	(SDV), in denen Fahrzeugfunktionen zunehmend als dynamisch gestartete Softwareservices
	realisiert werden. Die AUTOSAR Adaptive Platform (AP) unterstützt -- im Gegensatz zur
	AUTOSAR Classic Platform -- dynamisches Task-Management, Service-Discovery via SOME/IP
	und microservicebasierte Steuergerätearchitekturen. EDF$^+$ ist konzeptionell kompatibel
	mit dem AUTOSAR AP Execution Management, das Prozesse (\texttt{Adaptive Application})
	mit Scheduling-Policies verwaltet. Eine konkrete Integrationsroute ist die Implementierung
	von EDF$^+$ als custom Scheduling Policy im POSIX-konformen AP-OS-Interface
	(analog zu \texttt{SCHED\_DEADLINE} unter Linux oder QNX).
	Da ISO~26262 ASIL-D-Systeme eine vollständige Nachweisbarkeit der Scheduling-Latenz
	verlangt, bietet die in Satz~\ref{satz:hoeffding} hergeleitete Hoeffding-Schranke eine
	direkte Möglichkeit zur statistisch abgesicherten Worst-Case-Response-Time-Analyse
	(WCRT), wie sie im AUTOSAR Safety-Manifest dokumentiert werden muss.
	Die Selbststabilisierungseigenschaft (Satz~\ref{satz:selbst}) erleichtert darüber hinaus
	die Argumentation gegenüber dem Functional Safety Manager: Transiente Blocking-Ereignisse
	-- etwa durch Over-the-Air-Updates, dynamisch geladene WASM-Module oder kurzfristige
	Netzwerkverzögerungen im SOME/IP-Stack -- können das Scheduler-Verhalten nicht dauerhaft
	beeinträchtigen. Langfristig ist EDF$^+$ ein Kandidat für
	die Integration in den AUTOSAR AP Safety Manager als formell nachgewiesener, selbst\-stabilisierender Scheduler
	für Mixed-Criticality-Systeme (ASIL A bis D auf einem SoC).

	\paragraph{Formale Verifikation mit UPPAAL, SPIN und TLA$^+$.}
	Der Korrektheitsbeweis in Abschnitt~4 ist ein mathematischer Beweis auf Papier; für
	sicherheitskritische Systeme gemä\ss{} ISO~26262 ASIL-D wird zusätzlich ein
	maschinengeprüfter, formaler Nachweis gefordert. Drei Verifikationsansätze sind
	für EDF$^+$ besonders relevant:
	\begin{itemize}
		\item \textbf{UPPAAL (Timed Automata):} Der EDF$^+$-Algorithmus kann als Netzwerk
		      zeitbehafteter Automaten kodiert werden; die Tasks
		      $\{running, blocked, ready, finished\}$ entsprechen den Zuständen je eines
		      Automaten. Die Optimalitätseigenschaft ist als CTL-Formel
		      $\mathrm{AG}(\neg\, \textit{deadline\_miss})$ ausdrückbar.
		      Um der Zustandsraumexplosion bei $N = 40$ zu begegnen, bieten sich
		      parametrische Verifikation für kleine $N$ (etwa $N \le 8$) mit
		      anschlie\ss{}endem Induktionsschritt sowie abstrakte Interpretation an.
		\item \textbf{SPIN/Promela:} Blocking-Ereignisse können als nicht-deterministische
		      Transitionen in Promela kodiert werden. Die Fairness-Eigenschaft
		      (Satz~\ref{satz:fairness}) ist als LTL-Formel
		      $\mathrm{GF}(\textit{task\_executed})$ ausdrückbar und mit dem SPIN Model
		      Checker prüfbar. Besonders wertvoll ist SPIN für den Nachweis der
		      Deadlock-Freiheit der Neuplanungsschleife.
		\item \textbf{TLA$^+$ (Lamport):} TLA$^+$ bietet die höchste Ausdruckskraft für
		      verteilte und nebenläufige Algorithmen. EDF$^+$ lässt sich als
		      Zustandsmaschine mit explizitem Blocking-Modell vollständig in TLA$^+$
		      spezifizieren und mit dem TLC Model Checker prüfen.
	\end{itemize}
	Ein langfristiges Ziel ist die formale Verifikation im Proof-Assistant \emph{Coq} oder
	\emph{Isabelle/HOL}, was einen maschinengeprüften Korrektheitsbeweis liefern würde,
	der direkt in das Primärzertifizierungsdossier für ISO~26262 integrierbar ist.
	Die Kombination aus mathematischem Beweis (diese Arbeit) und Theorem-Prover-Zertifikat
	wäre ein Novum für Echtzeit-Scheduler in der Automobilindustrie.

	\paragraph{Approximative und ressourcenoptimierte Varianten für gro\ss{}e $N$.}
	Für Echtzeitsysteme mit sehr gro\ss{}em $N$ (etwa $N \ge 1000$ Tasks in
	Hochleistungssteuergeräten für autonomes Fahren) kann der $O(N \log N)$-Overhead
	jeder Neuplanung trotz asymptotischer Effizienz absolut betrachtet in die Nähe von
	Echtzeitleatenzgrenzen rücken. Folgende Optimierungen sind vielversprechend:
	\begin{itemize}
		\item \textbf{Inkrementelle Neuplanung:} Die Monotonieeigenschaft
		      (Abschnitt~5.3) besagt, dass nur die Priorität der einen
		      blockierten Task $T_i$ geändert wird. Statt Merge Sort vollständig zu
		      wiederholen, genügt das Einfügen von $T_i$ ans Ende einer bereits sortierten
		      Liste in $O(\log N)$ via binärer Suche -- eine Reduktion von $O(N \log N)$
		      auf $O(\log N)$ pro Blocking-Ereignis.
		\item \textbf{Prioritätsbasierte Datenstrukturen:} Min-Heaps oder
		      Fibonacci-Heaps unterstüt\-zen \texttt{insert}, \texttt{extract-min} und
		      \texttt{decrease-key} in amortisiert $O(\log N)$ bzw.\ $O(1)$ und können
		      die Warteschlange effizienter verwalten als ein sortiertes Array.
		\item \textbf{Probabilistische Zulassungskontrolle:} Für Systeme mit
		      hohem $p$ kann eine Admission-Control-Schicht implementiert werden: Tasks
		      werden nur zugelassen, wenn der mit Satz~\ref{satz:hoeffding} berechnete
		      WCRT-Bound die geforderte Schranke unterschreitet. Dies ermöglicht eine
		      lastadaptive Dämpfung von Overload-Situationen, wie sie bei dynamisch
		      zugeschalteten KI-Inferenz-Pipelines in autonomen Fahrzeugen auftreten können.
	\end{itemize}
	In Summe zeigt die Breite dieser Erweiterungsrichtungen, dass EDF$^+$ nicht nur ein
	akademischer Beitrag ist, sondern einen praktisch relevanten Ausgangspunkt für
	künftige Echtzeit-Scheduling-Forschung im Kontext eingebetteter, sicherheitskritischer
	und softwaregetriebener Fahrzeugsysteme darstellt.

	\paragraph{Weiterentwicklung von EDF$^+_{\mathrm{dyn}}$ für dynamische Task-Mengen.}
	Das in Kapitel~\ref{sec:dynamisch} eingeführte EDF$^+_{\mathrm{dyn}}$ eröffnet
	eine eigenständige, besonders reichhaltige Forschungsrichtung. Die dort bewiesene
	Optimalität (Satz~\ref{satz:optimal_dyn}) gilt unter der Annahme, dass ankommende
	Tasks mit erreichbaren Deadlines versehen sind und das System zu jedem Zeitpunkt
	schedulierbar bleibt ($U(t) \le 1$). Mehrere Erweiterungen dieser Grundaussage sind
	offen und praxisrelevant.

	\subparagraph{Overload-Management bei dynamisch wachsenden Task-Mengen.}
	In der Praxis kann die Systemauslastung durch unerwartet viele gleichzeitige
	Task-Ankünfte $U(t) > 1$ überschreiten. In diesem Fall ist es unmöglich, alle
	Deadlines einzuhalten. EDF$^+_{\mathrm{dyn}}$ bietet hier eine natürliche
	Anknüpfungsstelle für \emph{Overload-Management}-Strategien: Bei Prüfung der
	Schedulierbarkeit in Definition~\ref{def:sched_dyn} kann eine Admission-Control-Schicht
	neu ankommende Tasks ablehnen oder verzögern, falls die Aufnahme in $\mathcal{T}(t)$
	das System unschedulierbar machen würde. Formal lässt sich die Entscheidung durch
	den \emph{Acceptance-Test}
	\[
		U(t) + \frac{C_{\text{new}}}{D_{\text{new}} - t} \le 1
	\]
	algorithmisch in $O(1)$ prüfen (da $U(t)$ inkrementell gepflegt wird). Eine
	weitergehende Frage ist die \emph{optimale Verwerfungsstrategie} bei Overload:
	Welche Task soll abgewiesen werden, damit der kumulierte Nutzen aller rechtzeitig
	fertiggestellten Tasks maximiert wird? Dieser Ansatz verbindet EDF$^+_{\mathrm{dyn}}$
	mit der Theorie der \emph{Competitive Analysis} für Online-Algorithmen und
	liefert nützliche Gütegarantien selbst unter Worst-Case-Ankunftsfolgen.

	\subparagraph{Stochastische Analyse von Ankunftsprozessen.}
	Das Modell in Kapitel~\ref{sec:dynamisch} behandelt Task-Ankünfte als diskrete
	Ereignisse ohne stochastisches Modell. In der Praxis können Ankunftsprozesse durch
	Poisson-Prozesse ($\lambda$-parametrisiert), Markov-modulierte Prozesse (für
	bursthafte Ereignisfolgen in Kommunikationssystemen) oder erneuerungstheoretische
	Modelle (für sporadische Tasks mit bekanntem Mindestintervall $\tau_{\min}$)
	beschrieben werden. Die Verbindung mit der geometrischen Verteilung der
	Blocking-Ereignisse (Abschnitt~\ref{sec:dynamisch}) eröffnet die Möglichkeit,
	einen gemeinsamen stochastischen Rahmen für beide Ereignistypen zu entwickeln.
	Konkret könnte man den kombinierten Erwartungswert der Neuplanungen als
	\[
		E[C_{\mathrm{total}}^{\mathrm{dyn}}] = \bigl(1 + \lambda\mathcal{P} + Np\bigr)
		\cdot O(N_{\max}\log N_{\max})
	\]
	für einen Poisson-Ankunftsprozess mit Rate $\lambda$ ausdrücken, wobei
	$\lambda\mathcal{P}$ die erwartete Anzahl von Task-Ankünften pro Zyklus ist.
	Hoeffding-Bounds (Satz~\ref{satz:hoeffding}) könnten analog auf die Summe beider
	Ereignistypen ausgedehnt werden, um robuste WCRT-Abschätzungen zu ermöglichen.

	\subparagraph{Präemptionskosten und WCET-Modellierung.}
	In realen Prozessoren entstehen durch Präemption Overhead-Kosten: Cache-Räumung
	(Cache Flush), Pipeline-Invalidierung und Kontextsicherung. Das Modell in
	Algorithmus~\ref{alg:edfdyn} abstrahiert diese Kosten weg. Eine wichtige Erweiterung
	ist die Integration von \emph{Präemptionskosten} $\delta_{\mathrm{pre}} > 0$ in
	die WCET-Analyse: Die effektive Ausführungszeit einer Task $T_i$ erhöht sich
	durch jeden Präemptionsvorgang um $\delta_{\mathrm{pre}}$, was die tatsächliche
	Systemauslastung erhöht:
	\[
		U_{\mathrm{eff}}(t) = U(t) + k_{\mathrm{pre}}(t) \cdot \frac{\delta_{\mathrm{pre}}}{\mathcal{P}},
	\]
	wobei $k_{\mathrm{pre}}(t)$ die bisherige Anzahl von Präemptionen bezeichnet.
	Die Schedulierbarkeitsanalyse muss diese erhöhte Auslastung berücksichtigen;
	eine neue Schedulierbarkeitsgrenze $U_{\mathrm{max}}(p, \lambda, \delta_{\mathrm{pre}}) < 1$
	lässt sich analytisch herleiten und wäre ein wichtiges Ergebnis für den
	Einsatz in AUTOSAR ASIL-D-Systemen.

	\subparagraph{Hierarchisches EDF$^+_{\mathrm{dyn}}$ für Multi-Anwendungsszenarien.}
	Moderne SDV-Steuer\-geräte führen mehrere Anwendungen mit unterschiedlicher
	Kritikalität (ASIL A bis D) gleichzeitig aus. Eine hierarchische Erweiterung
	von EDF$^+_{\mathrm{dyn}}$ wäre wie folgt strukturiert: Auf der obersten Ebene
	kontrolliert ein \emph{Mixed-Criticality-Scheduler} (etwa nach Vestal
	\cite{Vestal2007} oder Burns \& Davis \cite{Burns2013}) die Zuweisung von
	CPU-Budgets an Anwendungspartitionen. Innerhalb jeder Partition läuft
	EDF$^+_{\mathrm{dyn}}$ als lokaler Scheduler und verwaltet die dynamisch
	ankommenden Tasks der Partition autonom. Neu ankommende Tasks von einem
	höheren Kritikalitätsniveau können die Partitionsgrenzen überschreiten
	und einen \emph{Cross-Partition-Präemptionsmechanismus} auslösen, der analog
	zur Ankunftspräemption in Definition~\ref{def:ankunft} formalisiert werden kann.
	Dieser Ansatz verbindet die bewiesene Optimalität von EDF$^+_{\mathrm{dyn}}$
	auf Partitionsebene mit der Zertifizierbarkeit gemischtkritischer Systeme nach
	ISO~26262 und DO-178C.

	\subparagraph{Formale Verifikation von EDF$^+_{\mathrm{dyn}}$.}
	Die Aussagen des dynamischen Modells -- insbesondere Satz~\ref{satz:optimal_dyn}
	und die Lemmata~\ref{lem:ankunft_konsistenz} und~\ref{lem:praemption_sicherheit}
	-- sind ebenso wie die statischen Sätze Kandidaten für eine formale Verifikation
	in UPPAAL, TLA$^+$ oder Isabelle/HOL. Die besondere Herausforderung liegt im
	unendlichen Zustandsraum, der durch beliebig viele Task-Ankünfte entsteht.
	Ein vielsprechender Ansatz sind \emph{parametrische zeitbehaftete Automaten},
	bei denen die Anzahl der Ankünfte $A$ als Parameter behandelt wird und
	Correctness-Eigenschaften als Funktionen von $A$ nachgewiesen werden können.
	Alternativ kann eine Abstraktion auf endliche Modelle (bounded model checking)
	mit wachsenden Schranken $A \le A_{\max}$ systematisch Gegenbeispiele ausschlie\ss{}en
	und so das Vertrauen in die Korrektheit des dynamischen Algorithmus sukzessive stärken.

	\subparagraph{Simulation und Benchmark-Evaluierung.}
	Die in Abschnitt~8 durchgeführten Monte-Carlo-Simulationen betreffen ausschlie\ss{}lich
	das statische Modell. Analoge Simulationen für EDF$^+_{\mathrm{dyn}}$ sollten
	für verschiedene Ankunftsraten $\lambda$, Blocking-Wahrscheinlichkeiten $p$ und
	Systemauslastungen $U$ durchgeführt werden. Besonders interessant sind
	Szenarien mit kombinierten Ankünfts- und Blocking-Bursts, wie sie in
	vernetzten Fahrzeugsystemen bei Feldbus-Ãœberlast oder OTA-Update-Kampagnen
	auftreten können. Ein Benchmark-Vergleich zwischen EDF$^+_{\mathrm{dyn}}$,
	dem klassischen dynamischen EDF sowie FIFO und Priority-Inheritance-basiertem
	RMS würde die theoretischen Optimalitätsaussagen empirisch absichern und
	die quantitative Ãœberlegenheit in praxisnahen Szenarien demonstrieren.

	\begin{thebibliography}{99}
		
		\bibitem{Liu1973}
		C.\,L.\ Liu und J.\,W.\ Layland,
		``Scheduling algorithms for multiprogramming in a hard-real-time environment,''
		\textit{Journal of the ACM}, Bd.~20, Nr.~1, S.~46--61, 1973.
		
		\bibitem{Buttazzo2011}
		G.\,C.\ Buttazzo,
		\textit{Hard Real-Time Computing Systems: Predictable Scheduling Algorithms and Applications},
		3.~Aufl.\ New York: Springer, 2011.
		
		\bibitem{Sha1990}
		L.\ Sha, R.\ Rajkumar und J.\,P.\ Lehoczky,
		``Priority inheritance protocols: An approach to real-time synchronization,''
		\textit{IEEE Transactions on Computers}, Bd.~39, Nr.~9, S.~1175--1185, 1990.
		
		\bibitem{Hoeffding1963}
		W.\ Hoeffding,
		``Probability inequalities for sums of bounded random variables,''
		\textit{Journal of the American Statistical Association}, Bd.~58, Nr.~301, S.~13--30, 1963.
		
		\bibitem{AUTOSAR2023}
		AUTOSAR Consortium,
		\textit{AUTOSAR Classic Platform -- Specification of Operating System},
		Release~R23-11, 2023.
		
		\bibitem{ISO26262}
		ISO,
		\textit{ISO~26262: Road Vehicles -- Functional Safety},
		2.~Aufl., 2018.
		
		\bibitem{Feller1968}
		W.\ Feller,
		\textit{An Introduction to Probability Theory and Its Applications},
		Bd.~1, 3.~Aufl.\ New York: Wiley, 1968.
		
		\bibitem{Cormen2022}
		T.\,H.\ Cormen, C.\,E.\ Leiserson, R.\,L.\ Rivest und C.\ Stein,
		\textit{Introduction to Algorithms},
		4.~Aufl.\ Cambridge, MA: MIT Press, 2022.

		\bibitem{Stankovic1988}
		J.\,A.\ Stankovic,
		``Misconceptions about real-time computing: A serious problem for next-generation systems,''
		\textit{IEEE Computer}, Bd.~21, Nr.~10, S.~10--19, 1988.

		\bibitem{Lehoczky1989}
		J.\,P.\ Lehoczky, L.\ Sha und Y.\ Ding,
		``The rate monotonic scheduling algorithm: Exact characterization and average case behavior,''
		in \textit{Proc.\ 10th IEEE Real-Time Systems Symposium (RTSS)}, S.~166--171, 1989.

		\bibitem{Dertouzos1974}
		M.\,L.\ Dertouzos,
		``Control robotics: The procedural control of physical processes,''
		in \textit{Proc.\ IFIP Congress}, S.~807--813, 1974.

		\bibitem{Spuri1996}
		M.\ Spuri und G.\ Buttazzo,
		``Scheduling aperiodic tasks in dynamic priority systems,''
		\textit{Real-Time Systems}, Bd.~10, Nr.~2, S.~179--210, 1996.

		\bibitem{Baker1991}
		T.\,P.\ Baker,
		``Stack-based scheduling of realtime processes,''
		\textit{Real-Time Systems}, Bd.~3, Nr.~1, S.~67--99, 1991.

		\bibitem{Abdelzaher2004}
		T.\ Abdelzaher, E.\,G.\ Coffman Jr.\ und J.\ Lehoczky,
		``Optimal scheduling of push and pull systems,''
		\textit{ACM SIGMETRICS Performance Evaluation Review}, Bd.~32, Nr.~1, S.~13--24, 2004.

		\bibitem{Audsley1993}
		N.\,C.\ Audsley, A.\ Burns, M.\ Richardson, K.\ Tindell und A.\,J.\ Wellings,
		``Applying new scheduling theory to static priority pre-emptive scheduling,''
		\textit{Software Engineering Journal}, Bd.~8, Nr.~5, S.~284--292, 1993.

		\bibitem{Baruah1990}
		S.\,K.\ Baruah, L.\,E.\ Rosier und R.\,R.\ Howell,
		``Algorithms and complexity concerning the preemptive scheduling of periodic,
		real-time tasks on one processor,''
		\textit{Real-Time Systems}, Bd.~2, Nr.~4, S.~301--324, 1990.

		\bibitem{Dhall1978}
		S.\,K.\ Dhall und C.\,L.\ Liu,
		``On a real-time scheduling problem,''
		\textit{Operations Research}, Bd.~26, Nr.~1, S.~127--140, 1978.

		\bibitem{Pellizzoni2008}
		R.\ Pellizzoni und G.\ Lipari,
		``Improved schedulability analysis of real-time transactions with earliest deadline scheduling,''
		\textit{ACM Transactions on Embedded Computing Systems}, Bd.~7, Nr.~2, S.~1--19, 2008.

		\bibitem{Brandenburg2011}
		B.\,B.\ Brandenburg und J.\,H.\ Anderson,
		``Optimality results for multiprocessor real-time locking,''
		in \textit{Proc.\ 32nd IEEE Real-Time Systems Symposium (RTSS)}, S.~49--60, 2011.

		\bibitem{Rajkumar1991}
		R.\ Rajkumar,
		\textit{Synchronization in Real-Time Systems: A Priority Inheritance Approach},
		Dordrecht: Kluwer Academic Publishers, 1991.

		\bibitem{Mok1983}
		A.\,K.\ Mok,
		\textit{Fundamental Design Problems of Distributed Systems for the Hard-Real-Time Environment},
		Diss., Massachusetts Institute of Technology, 1983.

		\bibitem{Bini2004}
		E.\ Bini und G.\,C.\ Buttazzo,
		``Schedulability analysis of periodic fixed priority systems,''
		\textit{IEEE Transactions on Computers}, Bd.~53, Nr.~11, S.~1462--1473, 2004.

		\bibitem{Anderson2005}
		J.\,H.\ Anderson, J.\,M.\ Calandrino und U.\,C.\ Devi,
		``Real-time scheduling on multicore platforms,''
		in \textit{Proc.\ 11th IEEE Real-Time and Embedded Technology and Applications Symposium (RTAS)},
		S.~179--190, 2005.

		\bibitem{Chetto1990}
		H.\ Chetto, M.\ Silly und T.\ Bouchentouf,
		``Dynamic scheduling of real-time tasks under precedence constraints,''
		\textit{Real-Time Systems}, Bd.~2, Nr.~3, S.~181--194, 1990.

		\bibitem{Harbour1991}
		M.\,G.\ Harbour, M.\,H.\ Klein und J.\,P.\ Lehoczky,
		``Timing analysis for fixed-priority scheduling of hard real-time systems,''
		\textit{IEEE Transactions on Software Engineering}, Bd.~20, Nr.~1, S.~13--28, 1994.

		\bibitem{Lipari2000}
		G.\ Lipari und S.\,K.\ Baruah,
		``A hierarchical extension to the earliest deadline first scheduling algorithm,''
		in \textit{Proc.\ 6th IEEE Real-Time Technology and Applications Symposium (RTAS)},
		S.~123--131, 2000.

		\bibitem{Baruah2011}
		S.\ Baruah, V.\ Bonifaci, G.\ D'Angelo, H.\ Li, A.\ Marchetti-Spaccamela,
		S.\ van der Ster und L.\ Stougie,
		``Preemptive uniprocessor scheduling of mixed-criticality sporadic task systems,''
		\textit{Journal of the ACM}, Bd.~62, Nr.~2, S.~1--33, 2015.

		\bibitem{Vestal2007}
		S.\ Vestal,
		``Preemptive scheduling of multi-criticality systems with varying degrees of
		execution time assurance,''
		in \textit{Proc.\ 28th IEEE Real-Time Systems Symposium (RTSS)}, S.~239--243, 2007.

		\bibitem{Burns2013}
		A.\ Burns und R.\ Davis,
		\textit{Mixed Criticality Systems -- A Review},
		v5.0, Univ.\ of York, Department of Computer Science, 2013.

	\end{thebibliography}
	
\end{document}
